\documentclass[
  aps,
  pra,
  reprint,
  superscriptaddress,
  amsmath,
  amssymb,
  nofootinbib,
  floatfix
]{revtex4-2}

\usepackage{amsthm}
\usepackage{bm}
\usepackage{booktabs}
\usepackage{tikz}
\usetikzlibrary{calc}

% Symmetric Yu--Oh figure, adapted from Svozil (2025), Fig. 1(a).
\newif\ifDistinctContexts
\DistinctContextsfalse

% Constant colours on context orbits preserve the threefold colour symmetry.
\definecolor{YOtriad}{HTML}{246A9A}
\definecolor{YObasis}{HTML}{8A939B}
\definecolor{YOspoke}{HTML}{237D70}
\definecolor{YOouterA}{HTML}{85458B}
\definecolor{YOouterB}{HTML}{B47620}
\definecolor{YOdiagonal}{HTML}{B44850}
\colorlet{YOctx1}{YOtriad}
\colorlet{YOctx2}{YOtriad}
\colorlet{YOctx3}{YOtriad}
\colorlet{YOctx4}{YObasis}
\colorlet{YOctx5}{YOspoke}
\colorlet{YOctx6}{YOspoke}
\colorlet{YOctx7}{YOspoke}
\colorlet{YOctx8}{YOouterA}
\colorlet{YOctx11}{YOouterA}
\colorlet{YOctx16}{YOouterA}
\colorlet{YOctx9}{YOouterB}
\colorlet{YOctx13}{YOouterB}
\colorlet{YOctx14}{YOouterB}
\colorlet{YOctx10}{YOdiagonal}
\colorlet{YOctx12}{YOdiagonal}
\colorlet{YOctx15}{YOdiagonal}
\ifDistinctContexts
  \definecolor{YOctx1}{HTML}{BF1FB4}
  \definecolor{YOctx2}{HTML}{B4000A}
  \definecolor{YOctx3}{HTML}{C99400}
  \definecolor{YOctx4}{HTML}{555F00}
  \definecolor{YOctx5}{HTML}{559F2A}
  \definecolor{YOctx6}{HTML}{D48A6A}
  \definecolor{YOctx7}{HTML}{F47400}
  \definecolor{YOctx8}{HTML}{B4944A}
  \definecolor{YOctx9}{HTML}{DF74BF}
  \definecolor{YOctx10}{HTML}{744ABF}
  \definecolor{YOctx11}{HTML}{35552A}
  \definecolor{YOctx12}{HTML}{C955FF}
  \definecolor{YOctx13}{HTML}{00B49F}
  \definecolor{YOctx14}{HTML}{6A558A}
  \definecolor{YOctx15}{HTML}{743F00}
  \definecolor{YOctx16}{HTML}{FF0035}
\fi
\tikzset{
  YOtrace/.style={line width=1.45pt,line cap=round,line join=round},
  YOapex/.style={line width=.68pt,line cap=round,dash pattern=on 2.4pt off 2.1pt,draw opacity=.65},
  YOlabel/.style={font=\small,inner sep=1.2pt,fill=white},
  YOnumber/.style={font=\sffamily\fontsize{8}{8}\selectfont,inner sep=0pt},
}
\newcommand{\YOline}[4]{\draw[YOtrace,YOctx#1] (v#2)--(v#3)--(v#4);}
\newcommand{\YOapexlink}[6]{%
  \draw[YOapex,draw=YOctx#1] (#3)
    -- ($(#3)!-#4!(#2)$)
    .. controls ($(#3)!-#5!(#2)$) and #6 .. (v26);%
}
\newcommand{\YOapexvia}[9]{%
  \draw[YOapex,draw=YOctx#1] (#3)
    -- ($(#3)!-#4!(#2)$)
    .. controls ($(#3)!-#5!(#2)$) and #6 .. #7
    .. controls #8 and #9 .. (v26);%
}
\newcommand{\YOvertex}[2]{%
  \foreach \context/\radius in {#2}{%
    \fill[YOctx\context] (v#1) circle[radius=\radius cm];%
  }%
  \fill[white] (v#1) circle[radius=.18cm];%
  \node[YOnumber] at (v#1) {#1};%
}


\usepackage{hyperref}

\hypersetup{
  colorlinks=true,
  linkcolor=blue,
  citecolor=blue,
  urlcolor=blue
}


\usepackage{orcidlink}

\newtheorem{theorem}{Theorem}
\newtheorem{proposition}[theorem]{Proposition}
\newtheorem{lemma}[theorem]{Lemma}
\newtheorem{corollary}[theorem]{Corollary}
\theoremstyle{definition}
\newtheorem{definition}[theorem]{Definition}
\theoremstyle{remark}
\newtheorem{remark}[theorem]{Remark}

\newcommand{\R}{\mathbb{R}}
\newcommand{\Q}{\mathbb{Q}}
\newcommand{\C}{\mathbb{C}}
\newcommand{\F}{\mathbb{F}}
\newcommand{\PP}{\mathbb{P}}
\newcommand{\YO}{\mathrm{YO}}
\DeclareMathOperator{\spanop}{span}
\DeclareMathOperator{\supp}{supp}

\begin{document}

\title{Four-Dimensional Context Hypergraphs with Separating Two-Valued States and a Chromatic Gap}

\author{Karl Svozil\,\orcidlink{0000-0001-6554-2802}}
\affiliation{Institute for Theoretical Physics, TU Wien,
Wiedner Hauptstra{\ss}e 8-10/136, A-1040 Vienna, Austria}
\email[Corresponding author: ]{karl.svozil@tuwien.ac.at}
\homepage{https://tph.tuwien.ac.at/~svozil}


\date{\today}

\begin{abstract}
We complete a known 17-ray graph into a linear hypergraph of 65 rays and 29 orthogonal bases in real dimension four. An integer realization has exactly the prescribed orthogonalities, and its graph has clique number four and chromatic number five. We determine its fractional chromatic number and all two-valued states; 48 states suffice for a unital, separating partition representation of the vertices and listed contexts. This representation fails order determination and identifies distinct compound Hilbert-space events. An apex suspension of the Yu--Oh configuration provides a smaller non-linear example with the same chromatic gap.
\end{abstract}

\maketitle

\section{Introduction and overview}

Let $H=(V,E)$, with $E\ne\emptyset$, be a $4$-uniform context hypergraph: each context $e\in E$ is a $4$-element subset of $V$, and its orthogonality graph $\Gamma(H)$ has vertex set $V$ and edge set
\begin{equation}\label{eq:orth-graph}
E(\Gamma(H))=\bigl\{\{u,v\}\subseteq V : u\ne v,\ \exists\,e\in E,\ \{u,v\}\subseteq e\bigr\}.
\end{equation}
Writing $K_t$ for the complete graph on $t$ vertices, every context $e\in E$ forms a $4$-clique ($K_4$) in $\Gamma(H)$, so the clique number satisfies $\omega(\Gamma(H))\ge 4$, and equality $\omega(\Gamma(H))=4$ holds whenever $\Gamma(H)$ is $K_5$-free. A \emph{strong $k$-coloring} of $H$ is a proper $k$-coloring $c:V\to\{0,1,\ldots,k-1\}$ of $\Gamma(H)$, so that every context $e\in E$ receives four distinct colors; the minimum such $k$ is the chromatic number $\chi(H)=\chi(\Gamma(H))\ge\omega(\Gamma(H))$.

In four-dimensional Hilbert space, a $4$-uniform hypergraph $H=(V,E)$ is \emph{edge-faithfully coordinatizable in $\R^4$} if there is an injective ray assignment $i\mapsto [v_i]\in\PP(\R^4)$, the projective space of rays in real Euclidean four-space ($v_i\in\R^4\setminus\{0\}$) such that
\begin{equation}\label{eq:edge-faithful}
v_i\cdot v_j = 0 \iff \{i,j\}\in E(\Gamma(H))
\qquad(i\ne j\in V).
\end{equation}
In particular, every context $e\in E$ is mapped to an orthogonal basis of $\R^4$, no two vertices share a ray, and no extra orthogonalities occur outside $E(\Gamma(H))$.

Independently of quantum coordinatization, we call $H$ \emph{set-representable on its vertices and listed contexts} if its collection $\mathcal{S}(H)$ of \emph{two-valued states}---maps $s:V\to\{0,1\}$ satisfying $\sum_{i\in e}s(i)=1$ for every $e\in E$---is both:
\begin{enumerate}
\item \emph{unital}: for every $i\in V$, there exists $s\in\mathcal{S}(H)$ with $s(i)=1$; and
\item \emph{separating}: for every pair of distinct vertices $i\ne j\in V$, there exists $s\in\mathcal{S}(H)$ with $s(i)\ne s(j)$.
\end{enumerate}
For any unital and separating family $\Omega\subseteq\mathcal{S}(H)$, set
\begin{equation}\label{eq:set-representation}
A_i=\{s\in\Omega:s(i)=1\}\qquad(i\in V).
\end{equation}
The sets $A_i$ are nonempty and pairwise distinct, and $\{A_i:i\in e\}$ is a partition of $\Omega$ for each $e\in E$. This is the precise partition representation meant here~\cite{svozil-tkadlec,svozil-2021-chroma}. It is weaker than an order-determining representation: inclusion $A_i\subseteq A_j$ can hold for distinct, hence incomparable, Hilbert-space rays. Both examples below fail order determination already on atoms; the linear example also identifies distinct compound Hilbert-space events. Section~\ref{sec:partition-scope} and Appendix~\ref{app:states-hyo4} give explicit witnesses. Their separating states therefore do not supply a faithful Boolean representation of the Hilbert-space event structure or a model of all quantum preparations. The support of a state is $\supp(s)=\{i\in V:s(i)=1\}$.

The \emph{independence number} $\alpha(G)$ is the largest size of a vertex set containing no adjacent pair. The \emph{fractional chromatic number} $\chi_f(G)$ is the minimum total nonnegative weight assigned to independent sets so that the sets containing each vertex have total weight at least one.

Because the color classes of any strong $4$-coloring of a $4$-uniform hypergraph define four two-valued states that partition $V$, chromatic completeness ($\chi(H)=4$) implies unitality, whereas Kochen--Specker (KS) configurations such as the $24$-ray Peres set in $\R^4$~\cite{peres-91} satisfy $\chi>4$ because $\mathcal{S}(H)=\emptyset$. The subscript $\YO$ abbreviates Yu--Oh. Conversely, the existence of unital and separating states need not supply a strong $4$-coloring: such a coloring requires four states whose supports are pairwise disjoint and cover $V$. When applied to the completion or the selected state family below, \emph{canonical} refers to the specified coordinate and indexing conventions, rather than to an isomorphism-invariant choice.

\begin{table*}[t]
\caption{Comparison of $4$-uniform context hypergraphs in $\R^4$ with clique number $\omega(\Gamma(H))=4$ and chromatic number $\chi(H)=5>4$, showing the number of pairwise non-collinear vectors $|V|$, the number of $4$-element contexts $|E|$, and the number of orthogonality edges $|E(\Gamma(H))|$.}\label{tab:comparison}
\begin{ruledtabular}
\scriptsize
\begin{tabular}{lccccccccc}
Hypergraph & $|V|$ & $|E|$ & $|E(\Gamma(H))|$ & Linear & $\omega(\Gamma(H))$ & $\chi(H)$ & Edge-faithful & $|\mathcal{S}(H)|$ & Set-representable \\ \hline
Peres $24$-ray set~\cite{peres-91} & $24$ & $24$ & $108$ & No ($|e\cap e'|\le2$) & $4$ & $5$ & Yes ($\R^4$) & $0$ & No (KS) \\
Apex suspension $H_{\YO}^{(4)}$ (Prop.~\ref{prop:warmup}) & $26$ & $16$ & $73$ & No ($|e\cap e'|\le2$) & $4$ & $5$ & Yes ($\R^4$) & $25$ & Yes \\
Canonical linear $H_{65}$ (Thm.~\ref{thm:h65}) & $65$ & $29$ & $174$ & Yes ($|e\cap e'|\le1$) & $4$ & $5$ & Yes ($\R^4$) & $393{,}216$ & Yes ($48$ canonical)
\end{tabular}
\end{ruledtabular}
\end{table*}

The purpose of this note is to exhibit explicit $4$-uniform context hypergraphs $H$ with $\omega(\Gamma(H))=4$ that are simultaneously set-representable (unital and separating two-valued states), edge-faithfully coordinatizable in $\R^4$, and chromatically incomplete ($\chi(H)=5>4=\omega(\Gamma(H))$). Table~\ref{tab:comparison} summarizes the configurations analyzed below.

The chromatic obstructions are inherited from the $17$-ray graph of Haynes \emph{et al.}~\cite{haynes2010orthogonal} and the Yu--Oh configuration~\cite{Yu-2012}; related distinctions between states and colorings appear in Refs.~\cite{svozil-2021-chroma,svozil-2025-color}. The additional assertions proved here concern linear context completion: a rational two-plane lemma supplies integer vectors preserving exact orthogonalities and chromatic colorability, with a criterion for extending separating states. The $65$-ray example realizes these properties with explicit state and fractional-coloring certificates.

\section{Non-linear warm-up \texorpdfstring{$H_{\YO}^{(4)}$}{H\_YO\^(4)} and the hyperplane obstruction}

Before constructing a linear $4$-uniform example, we record an elementary $26$-ray, $16$-context non-linear example obtained by lifting the $3$-uniform completed Yu--Oh hypergraph $H_{\YO}^{+}=(V_{\YO}^{+},E_{\YO}^{+})$~\cite{Yu-2012,Pavi_i__2019,svozil-2025-color} from $\R^3$ to $\R^4$. The obstruction below rules out retaining linearity when two intersecting complete orthogonal triads in one hyperplane are lifted to distinct four-element contexts.

\begin{proposition}[Non-linear apex suspension $H_{\YO}^{(4)}$]\label{prop:warmup}
Let $H_{\YO}^{+}=(V_{\YO}^{+},E_{\YO}^{+})$ be the $25$-vertex, $16$-context completed Yu--Oh hypergraph in $\R^3$, with vectors $r_1,\ldots,r_{25}\in\R^3$ and contexts $E_{\YO}^{+}=\{e_1,\ldots,e_{16}\}$. Define $H_{\YO}^{(4)}=(V_{26},E_{16}^{(4)})$ on $V_{26}=\{1,\ldots,26\}$ by appending an apex vertex $26$ to every context:
\begin{equation}\label{eq:hyo4-def}
E_{16}^{(4)} = \bigl\{e\cup\{26\} : e\in E_{\YO}^{+}\bigr\},
\end{equation}
and assign the vectors $v_i=(r_i,0)^T\in\R^4$ for $1\le i\le 25$ and $v_{26}=(0,0,0,1)^T\in\R^4$ (enumerated explicitly in Appendix~\ref{app:hyo4} and depicted as a $4$-uniform hypergraph diagram in Fig.~\ref{fig:hyo4}). Then:
\begin{enumerate}
\item $H_{\YO}^{(4)}$ is $4$-uniform with clique number $\omega(\Gamma(H_{\YO}^{(4)}))=4$ and chromatic number $\chi(H_{\YO}^{(4)})=\chi(H_{\YO}^{+})+1=5>4$.
\item The vectors $v_1,\ldots,v_{26}\in\R^4$ form an injective, edge-faithful orthogonal representation of $H_{\YO}^{(4)}$ in $\R^4$.
\item $H_{\YO}^{(4)}$ admits exactly $25$ two-valued states (enumerated in Appendix~\ref{app:states-hyo4}), which form a unital and separating collection on $V_{26}$.
\end{enumerate}
\end{proposition}

\begin{proof}
(1)--(2) Since $v_{26}=(0,0,0,1)^T$ is orthogonal to $\R^3\times\{0\}$ and $r_1,\ldots,r_{25}$ is an injective, edge-faithful representation of $H_{\YO}^{+}$ in $\R^3$, the vectors $v_1,\ldots,v_{26}$ are pairwise non-proportional and satisfy $v_i\cdot v_j=0$ if and only if $\{i,j\}\in E(\Gamma(H_{\YO}^{(4)}))$. Thus $\Gamma(H_{\YO}^{(4)})$ is the graph join $\Gamma(H_{\YO}^{+})*K_1$, which immediately gives $\omega(\Gamma(H_{\YO}^{(4)}))=\omega(\Gamma(H_{\YO}^{+}))+1=3+1=4$ and $\chi(H_{\YO}^{(4)})=\chi(H_{\YO}^{+})+1=4+1=5$.

(3) In any two-valued state $s:V_{26}\to\{0,1\}$ on $H_{\YO}^{(4)}$, either $s(26)=1$, which forces $s(1)=\cdots=s(25)=0$ (giving one apex state $s_{\mathrm{apex}}$), or $s(26)=0$, in which case $\sum_{i\in e}s(i)=1$ for all $e\in E_{\YO}^{+}$ so that $s|_{V_{\YO}^{+}}$ is one of the $24$ two-valued states of $H_{\YO}^{+}$. Since the $24$ states of $H_{\YO}^{+}$ are unital and separating on $\{1,\ldots,25\}$ and all vanish at $26$, while $s_{\mathrm{apex}}(26)=1$ and $s_{\mathrm{apex}}(i)=0$ for $i\le 25$, the $25$ states of $H_{\YO}^{(4)}$ are unital and separating on all $26$ vertices.
\end{proof}

Although $H_{\YO}^{(4)}$ satisfies all three requirements of set-representability, edge-faithfulness in $\R^4$, and $\chi=5>4=\omega$, it is \emph{non-linear}: the apex vertex $26$ lies in all $16$ contexts, so $|(e\cup\{26\})\cap(e'\cup\{26\})|=2$ whenever $e,e'\in E_{\YO}^{+}$ share a vertex in $\R^3$ (giving Berge girth $2$). A Berge cycle of length $r$ consists of $r$ distinct contexts and $r$ distinct vertices, with successive contexts meeting at the corresponding vertices cyclically; the Berge girth is the shortest such length. The following proposition shows that this non-linearity is unavoidable whenever two contexts have intersecting three-dimensional restrictions inside a common hyperplane.

\begin{proposition}[Hyperplane obstruction to linearity in dimension four]\label{prop:hyperplane-obstruction}
Let $H=(V,E)$ be a $4$-uniform context hypergraph with an injective orthogonal representation $i\mapsto[v_i]\in\PP(\F^4)$ ($\F\in\{\R,\C\}$, with $\C$ the complex numbers). Orthogonality over $\C$ is defined by the Hermitian product $\langle x,y\rangle=\sum_{j=1}^4\overline{x_j}y_j$. For a subspace $W$, write $e_W=\{i\in e:v_i\in W\}$. If two distinct contexts $e,e'\in E$ each contain at least three vertices whose rays lie in a common three-dimensional hyperplane $W=z^\perp\subset\F^4$ ($z\in\F^4\setminus\{0\}$), then both $e$ and $e'$ must contain the unique ray $[z]$ orthogonal to $W$. Consequently, if $|e_W\cap e'_W|\ge1$, then $|e\cap e'|\ge2$, violating linearity.
\end{proposition}

\begin{proof}
Let $e=\{u_1,u_2,u_3,u_4\}$ with $v_{u_1},v_{u_2},v_{u_3}\in W$. Because $v_{u_1},v_{u_2},v_{u_3}$ are mutually orthogonal nonzero vectors in the $3$-dimensional subspace $W$, they form an orthogonal basis of $W$. Since $v_{u_4}\perp v_{u_1},v_{u_2},v_{u_3}$, we have $v_{u_4}\in W^\perp=\spanop(z)$, so $[v_{u_4}]=[z]$. The same argument applied to $e'$ shows that its fourth ray must also equal $[z]$; by injectivity, the same vertex $z_0\in V$ with $[v_{z_0}]=[z]$ belongs to both $e$ and $e'$. If $e$ and $e'$ also share any vertex inside $W$, then $|e\cap e'|\ge2$.
\end{proof}

\section{Construction of the canonical linear hypergraph \texorpdfstring{$H_{65}$}{H\_65} in \texorpdfstring{$\R^4$}{R4}}

By Proposition~\ref{prop:hyperplane-obstruction}, two intersecting complete orthogonal triads in one hyperplane cannot be extended to distinct contexts of a linear $4$-uniform hypergraph with an injective orthogonal representation in $\R^4$. The assertion concerns triads, not arbitrary smaller restrictions of contexts. Instead, we start from a four-dimensional core of $17$ rays in $\{-1,0,1\}^4$ and complete its $24$ uncovered orthogonal pairs into $4$-contexts via the following general lemma. All completion vectors are integer vectors, with $\mathbb{Z}$ denoting the integers.

\subsection{The two-plane completion lemma in \texorpdfstring{$\R^4$}{R4}}

In $\R^3$, two orthogonal rays $\{v_u,v_v\}$ uniquely determine their third orthogonal basis ray $[v_u\times v_v]$, which may inadvertently be orthogonal or parallel to other rays of the configuration. In $\R^4$, by contrast, the orthogonal complement $\{v_u,v_v\}^\perp$ of an orthogonal pair is a \emph{two-dimensional plane}, providing an infinite one-parameter family of orthogonal completion pairs $\{w_1,w_2\}$.

\begin{lemma}[Simultaneous plane-completion faithfulness and state extension in $\R^4$]\label{lem:plane-completion}
Let $V_{\mathrm{core}}$ be a finite vertex set carrying pairwise non-collinear vectors $v_i\in\Q^4\setminus\{0\}$, where $\Q$ denotes the rational numbers, let $E_{\mathrm{base}}$ be a collection of $4$-element orthogonal bases of $\R^4$ drawn from $V_{\mathrm{core}}$ with $|C\cap C'|\le1$ for all $C\ne C'\in E_{\mathrm{base}}$, and let
\begin{equation}\label{eq:rem-edges}
E_{\mathrm{bridge}} = \bigl\{\{u,v\}\subset V_{\mathrm{core}} : v_u\cdot v_v=0\bigr\}\setminus \bigcup_{C\in E_{\mathrm{base}}}\binom{C}{2}
\end{equation}
be the set of orthogonal pairs in $V_{\mathrm{core}}$ not already covered by $E_{\mathrm{base}}$, with $E_{\mathrm{bridge}}=\{e_1,\ldots,e_M\}$ ($e_m=\{u_m,v_m\}$). For each $m\in\{1,\ldots,M\}$, set $L_m=\spanop(v_{u_m},v_{v_m})$ and $P_m=L_m^\perp$.
\begin{enumerate}
\item \emph{Span conditions for simultaneous edge-faithful completion:} The pairs $e_1,\ldots,e_M\in E_{\mathrm{bridge}}$ can be simultaneously completed to $4$-element orthogonal bases $E_m=\{u_m,v_m,w_{m,1},w_{m,2}\}$ by fresh integer vectors $w_{m,1},w_{m,2}\in P_m\cap\mathbb{Z}^4$ having \emph{no} extra orthogonalities or collinearities with any other vectors in the completed hypergraph $H=(V_{\mathrm{core}}\cup\bigcup_{m=1}^M\{w_{m,1},w_{m,2}\},\ E_{\mathrm{base}}\cup\{E_1,\ldots,E_M\})$ if and only if both
\begin{equation}\label{eq:span-condition}
v_z\notin L_m\quad\text{for every }z\in V_{\mathrm{core}}\setminus\{u_m,v_m\}
\quad(1\le m\le M)
\end{equation}
and
\begin{equation}\label{eq:plane-nonorth}
P_\ell\ne L_m\ (\text{i.e., }L_\ell\not\perp L_m)
\quad(1\le\ell\ne m\le M).
\end{equation}
\item \emph{Linearity, clique number, and chromatic preservation:} The completed $4$-uniform hypergraph $H$ is linear ($|e\cap e'|\le1$ for all $e\ne e'\in E(H)$). If $E(H)\ne\emptyset$ and the core orthogonality graph $G_{\mathrm{core}}$ is $K_5$-free, then $\omega(\Gamma(H))=4$, and for every $k\ge4$, $H$ is strongly $k$-colorable if and only if $G_{\mathrm{core}}$ is properly $k$-colorable.
\item \emph{Canonical state extension:} Let $S_1,\ldots,S_N\subset V_{\mathrm{core}}$ be all independent sets in $G_{\mathrm{core}}$ satisfying $|S_\alpha\cap C|=1$ for every $C\in E_{\mathrm{base}}$. For each $\alpha\in\{1,\ldots,N\}$ and $\varepsilon\in\{0,1\}$, define $s_{\alpha,\varepsilon}:V(H)\to\{0,1\}$ by $s_{\alpha,\varepsilon}(i)=\mathbf{1}_{i\in S_\alpha}$ for $i\in V_{\mathrm{core}}$ and, for each $m\in\{1,\ldots,M\}$,
\begin{equation}\label{eq:canonical-states}
\begin{aligned}
\bigl(s_{\alpha,\varepsilon}&(w_{m,1}),\,s_{\alpha,\varepsilon}(w_{m,2})\bigr)\\
&=\begin{cases}
(0,0), & |S_\alpha\cap\{u_m,v_m\}|=1,\\
(1-\varepsilon,\,\varepsilon), & S_\alpha\cap\{u_m,v_m\}=\emptyset.
\end{cases}
\end{aligned}
\end{equation}
Then every $s_{\alpha,\varepsilon}$ is a valid two-valued state on $H$; here $\mathbf{1}_{i\in S_\alpha}$ is the indicator of membership in $S_\alpha$. The two indices $\varepsilon=0,1$ can give the same state when $S_\alpha$ hits every bridge edge. If $S_1,\ldots,S_N$ are unital and separating on $V_{\mathrm{core}}$, and the $M$ edge zero-profiles
\begin{equation}\label{eq:zero-profiles}
Z(e_m) = \bigl\{\alpha\in\{1,\ldots,N\} : u_m\notin S_\alpha\text{ and }v_m\notin S_\alpha\bigr\}
\end{equation}
are nonempty and pairwise distinct for $m=1,\ldots,M$, then the indexed canonical family $\{s_{\alpha,\varepsilon}\}$, containing at most $2N$ distinct states, is unital and separating on all $|V_{\mathrm{core}}|+2M$ vertices of $H$.
\end{enumerate}
\end{lemma}

\begin{proof}
(1) For necessity, if $v_z\in L_m$ for some $z\in V_{\mathrm{core}}\setminus\{u_m,v_m\}$, then $P_m=L_m^\perp\subset v_z^\perp$, so every completion vector in $P_m$ is orthogonal to $v_z$; similarly, if $P_\ell=L_m$ for $\ell\ne m$, then $P_\ell\perp P_m$, forcing $w_{\ell,a}\perp w_{m,b}$ for all $a,b\in\{1,2\}$.

For sufficiency, assume Eqs.~\eqref{eq:span-condition}--\eqref{eq:plane-nonorth}. Because $v_{u_m},v_{v_m}\in\Q^4$, each $P_m$ is a rational $2$-plane and therefore admits an integer orthogonal basis $\{\eta_A,\eta_B\}\subset P_m\cap\mathbb{Z}^4$ (obtained by rational Gram--Schmidt and clearing denominators) with $n_A=\|\eta_A\|^2$ and $n_B=\|\eta_B\|^2$. We choose $\{w_{m,1},w_{m,2}\}\subset P_m\cap\mathbb{Z}^4$ inductively for $m=1,\ldots,M$ so as to maintain the invariant that, for every $r\ge m$, no previously placed vector $z\notin\{v_{u_r},v_{v_r}\}$ lies in $L_r$ (which holds at $m=1$ by Eq.~\eqref{eq:span-condition}). At step $m$, the invariant for $r=m$ ensures that every previously placed vector $z\notin\{v_{u_m},v_{v_m}\}$ satisfies $z\notin L_m=P_m^\perp$, so $z^\perp\cap P_m$ and $\R z\cap P_m$ are at most $1$-dimensional lines in $P_m$. Moreover, for every future $r>m$, Eq.~\eqref{eq:plane-nonorth} gives $P_m\ne L_r$, so $P_m\cap L_r$ is also at most a $1$-dimensional line in $P_m$. The integer parameterization
\begin{equation}\label{eq:plane-param}
\begin{aligned}
w_{m,1}(a,b)&=a\,\eta_A+b\,\eta_B,\\
w_{m,2}(a,b)&=-n_B b\,\eta_A+n_A a\,\eta_B
\end{aligned}
\end{equation}
produces infinitely many projectively distinct orthogonal pairs in $P_m\cap\mathbb{Z}^4$ as $(a,b)$ ranges over $\mathbb{Z}^2\setminus\{(0,0)\}$, and can avoid every forbidden line for both vectors: the inverse image of each such line under either invertible linear map in Eq.~\eqref{eq:plane-param} is a line in the $(a,b)$-parameter plane, and finitely many lines cannot cover all integer slopes. This avoids the finite union of forbidden lines in $P_m$, preventing any extra orthogonalities or collinearities with previously placed vectors while preserving $w_{m,1},w_{m,2}\notin L_r$ for all $r>m$.

(2) Each $E_m$ shares at most one vertex with any $C\in E_{\mathrm{base}}$ (since $\binom{C}{2}\cap E_{\mathrm{bridge}}=\emptyset$) and at most one vertex with any $E_{m'}$ ($m'\ne m$), while $w_{m,1},w_{m,2}$ have degree $1$ in $H$. In any proper $k$-coloring ($k\ge4$) of $G_{\mathrm{core}}$, $c(u_m)\ne c(v_m)$, leaving $k-2\ge2$ colors in $\{0,\ldots,k-1\}\setminus\{c(u_m),c(v_m)\}$ to color $w_{m,1}$ and $w_{m,2}$.

(3) Since $S_\alpha$ is independent in $G_{\mathrm{core}}$, $|S_\alpha\cap\{u_m,v_m\}|\in\{0,1\}$, so Eq.~\eqref{eq:canonical-states} gives $s_{\alpha,\varepsilon}(u_m)+s_{\alpha,\varepsilon}(v_m)+s_{\alpha,\varepsilon}(w_{m,1})+s_{\alpha,\varepsilon}(w_{m,2})=1$ on every $E_m$. Across the $2N$ indexed states $(s_{1,0},s_{1,1},\ldots,s_{N,0},s_{N,1})$, the profile of $w_{m,1}$ has $(1,0)$ in block $\alpha$ for $\alpha\in Z(e_m)$ and $(0,0)$ otherwise; the profile of $w_{m,2}$ has $(0,1)$ in block $\alpha$ for $\alpha\in Z(e_m)$ and $(0,0)$ otherwise; and the profile of $i\in V_{\mathrm{core}}$ has $(1,1)$ for $i\in S_\alpha$ and $(0,0)$ otherwise. Since $Z(e_m)\ne\emptyset$ and $Z(e_m)\ne Z(e_{m'})$ for $m\ne m'$, all $|V_{\mathrm{core}}|+2M$ state profiles are nonzero and pairwise distinct.
\end{proof}

\subsection{Explicit definition of \texorpdfstring{$H_{65}$}{H\_65} and provenance of the \texorpdfstring{$17$}{17}-ray core}

We define the $17$-vertex core $V_{\mathrm{core}}=\{1,\ldots,17\}$ in $\R^4$ using four structural groups of vectors in $\{-1,0,1\}^4$:
\begin{enumerate}
\item \emph{Standard coordinate basis $C_1=\{1,2,3,4\}$:}
\begin{align*}
v_1&=(1,0,0,0),& v_2&=(0,1,0,0),\\
v_3&=(0,0,1,0),& v_4&=(0,0,0,1).
\end{align*}
\item \emph{Even-parity Hadamard basis $C_2=\{5,6,7,8\}$} (the ray representatives with first coordinate $+1$ and an even number of minus signs):
\begin{align*}
v_5&=(1,1,1,1),& v_6&=(1,1,-1,-1),\\
v_7&=(1,-1,1,-1),& v_8&=(1,-1,-1,1).
\end{align*}
\item \emph{Odd-parity Hadamard hub triad $\{9,10,11\}$} (three mutually orthogonal $\pm1$ vectors with an odd number of minus signs):
\begin{align*}
v_9&=(1,1,-1,1),& v_{10}&=(1,-1,1,1),\\
v_{11}&=(1,-1,-1,-1).
\end{align*}
\item \emph{Three binary bridge pairs $B_1=\{12,13\}$, $B_2=\{14,15\}$, $B_3=\{16,17\}$} completing the three pairs of $\{9,10,11\}$ into $4$-element orthogonal bases $C_3=\{9,10,12,13\}$, $C_4=\{9,11,14,15\}$, and $C_5=\{10,11,16,17\}$:
\begin{align*}
v_{12}&=(1,0,0,-1),& v_{13}&=(0,1,1,0),\\
v_{14}&=(1,0,1,0),& v_{15}&=(0,1,0,-1),\\
v_{16}&=(1,1,0,0),& v_{17}&=(0,0,1,-1).
\end{align*}
\end{enumerate}

As a $17$-ray orthogonality graph in $\R^4$, $G_{\mathrm{core}}$ is orthogonally equivalent to the $17$-ray subgraph of the $24$-ray Peres configuration introduced by Haynes \emph{et al.}~\cite[Sec.~1.2, pp.~8--9]{haynes2010orthogonal} as an example with orthogonal dimension $4$ and chromatic number $5$. Explicitly, the orthogonal matrix
\begin{equation}\label{eq:haynes-transform}
T = \frac{1}{2}\begin{pmatrix}
 1 &  1 &  1 & -1\\
 1 &  1 & -1 &  1\\
 1 & -1 &  1 &  1\\
 1 & -1 & -1 & -1
\end{pmatrix},\qquad T^T T = I_4,
\end{equation}
where $I_4$ is the $4\times4$ identity matrix. The matrix $T$ maps the hub triad $\{v_9,v_{10},v_{11}\}$ to the coordinate rays $\{e_2,e_3,e_4\}$, the bases $C_1\cup C_2=\{v_1,\ldots,v_8\}$ to the eight rays $(1,\pm1,\pm1,\pm1)^T$, and the six bridge rays $\{v_{12},\ldots,v_{17}\}$ to the six binary rays $e_1\pm e_j$ ($j=2,3,4$) of Ref.~\cite{haynes2010orthogonal}. Choosing the coordinate-Hadamard frame above makes the bipartite parity structure between $C_1$, $C_2$, and $B_1\cup B_2\cup B_3$ transparent, allowing us to complete $G_{\mathrm{core}}$ into the linear $4$-uniform context hypergraph $H_{65}$ and analyze its two-valued state space by hand.

\begin{widetext}
\begin{equation}\label{eq:h65-coordinates}
\begin{gathered}
\begin{array}{c|rrrrrrrrrrrrrrrrr}
 i    & 1& 2& 3& 4& 5& 6& 7& 8& 9&10&11&12&13&14&15&16&17\\ \hline
 x_{i,1} & 1& 0& 0& 0& 1& 1& 1& 1& 1& 1& 1& 1& 0& 1& 0& 1& 0\\
 x_{i,2} & 0& 1& 0& 0& 1& 1&-1&-1& 1&-1&-1& 0& 1& 0& 1& 1& 0\\
 x_{i,3} & 0& 0& 1& 0& 1&-1& 1&-1&-1& 1&-1& 0& 1& 1& 0& 0& 1\\
 x_{i,4} & 0& 0& 0& 1& 1&-1&-1& 1& 1& 1&-1&-1& 0& 0&-1& 0&-1\\ \hline
 c(i) & 0& 1& 2& 3& 0& 1& 2& 4& 0& 1& 3& 3& 4& 2& 1& 0& 2
\end{array}
\\[1.2ex]
\begin{array}{c|rrrrrrrrrrrrrrrrrrrrrrrr}
 i    &18&19&20&21&22&23&24&25&26&27&28&29&30&31&32&33&34&35&36&37&38&39&40&41\\ \hline
 x_{i,1} & 1& 3& 1& 3& 0& 0& 2& 3& 1& 3& 1& 3& 0& 0& 2& 3& 1& 3& 1& 3& 0& 0& 2& 3\\
 x_{i,2} & 0& 0& 3&-2& 1& 3&-3& 1& 0& 0&-3& 2& 1& 3& 3&-1&-1&-3&-1&-3& 2& 3& 0& 0\\
 x_{i,3} &-3& 2& 0& 0&-1&-3& 3&-1&-1&-3&-1&-3& 3&-2& 0& 0& 0& 0& 3&-2& 3&-1&-3& 1\\
 x_{i,4} & 1& 3& 1& 3& 3&-2& 0& 0&-3& 2& 0& 0& 1& 3& 3&-1& 3&-2& 0& 0& 3&-1&-3& 1\\ \hline
 c(i) & 0& 2& 0& 1& 1& 2& 0& 1& 0& 3& 0& 1& 2& 3& 0& 3& 1& 3& 1& 2& 1& 3& 0& 3
\end{array}
\\[1.2ex]
\begin{array}{c|rrrrrrrrrrrrrrrrrrrrrrrr}
 i    &42&43&44&45&46&47&48&49&50&51&52&53&54&55&56&57&58&59&60&61&62&63&64&65\\ \hline
 x_{i,1} & 1& 2& 1& 2& 1& 5& 3& 3& 1& 2& 1& 2& 3& 3& 1& 5& 1& 2& 1& 2& 1& 5& 3& 3\\
 x_{i,2} &-5&-1&-3& 3&-2&-1&-1& 2& 3&-3& 5& 1&-2& 1& 2& 1&-1&-2&-1&-2& 3&-3&-5& 1\\
 x_{i,3} & 3&-3& 5& 1& 2& 1& 1&-2&-1&-2&-1&-2& 1&-5& 3&-3&-5&-1&-3& 3&-2&-1&-1& 2\\
 x_{i,4} & 1& 2& 1& 2&-3& 3& 5&-1& 5& 1& 3&-3&-2& 1& 2& 1&-3& 3&-5&-1&-2&-1&-1& 2\\ \hline
 c(i) & 1& 2& 0& 1& 0& 2& 0& 1& 0& 3& 0& 1& 2& 3& 0& 3& 1& 3& 1& 2& 1& 3& 0& 3
\end{array}
\end{gathered}
\end{equation}
\end{widetext}

Here the weight of a vector is its number of nonzero coordinates, and $N(i)=\{j\in V_{\mathrm{core}}\setminus\{i\}:v_i\cdot v_j=0\}$ denotes its neighborhood in $G_{\mathrm{core}}$. Because every vector in $C_1$ has weight $1$ and every vector in $C_2\cup\{9,10,11\}$ has weight $4$ with odd/even parity, there are no orthogonalities among $\{1,\ldots,11\}$ outside $C_1$, $C_2$, and the triad $\{9,10,11\}$. Each of the six binary bridge vectors $b\in\{12,\ldots,17\}$, having weight $2$, is orthogonal to exactly two coordinate rays in $C_1$ and two even-parity Hadamard rays in $C_2$, creating $6\times 2+6\times 2=24$ bridge edges $E_{\mathrm{bridge}}=\{e_1,\ldots,e_{24}\}$. Completing each $e_m=\{u_m,v_m\}$ via Lemma~\ref{lem:plane-completion} with fresh integer rays $\{2m+16,\,2m+17\}\subset\{18,\ldots,65\}$ yields the $29$ contexts of $H_{65}=(V_{65},E_{29})$:
\begin{equation}\label{eq:h65-contexts}
\begin{aligned}
E_{29}=\bigl\{
&\{1,2,3,4\},\ \{5,6,7,8\},\ \{9,10,12,13\},\\
&\{9,11,14,15\},\ \{10,11,16,17\},\ \{2,12,18,19\},\\
&\{3,12,20,21\},\ \{1,13,22,23\},\ \{4,13,24,25\},\\
&\{2,14,26,27\},\ \{4,14,28,29\},\ \{1,15,30,31\},\\
&\{3,15,32,33\},\ \{3,16,34,35\},\ \{4,16,36,37\},\\
&\{1,17,38,39\},\ \{2,17,40,41\},\ \{5,12,42,43\},\\
&\{8,12,44,45\},\ \{6,13,46,47\},\ \{7,13,48,49\},\\
&\{6,14,50,51\},\ \{8,14,52,53\},\ \{5,15,54,55\},\\
&\{7,15,56,57\},\ \{7,16,58,59\},\ \{8,16,60,61\},\\
&\{5,17,62,63\},\ \{6,17,64,65\}
\bigr\}.
\end{aligned}
\end{equation}
The complete coordinate table $v_i=(x_{i,1},x_{i,2},x_{i,3},x_{i,4})^T\in\R^4$ and a strong $5$-coloring $c(i)\in\{0,1,2,3,4\}$ for all $65$ vertices are given in Eq.~\eqref{eq:h65-coordinates} (and enumerated vertex-by-vertex and context-by-context in Appendix~\ref{app:h65}). Figure~\ref{fig:hyo4} illustrates the $26$-vertex, $16$-context warm-up apex suspension $H_{\YO}^{(4)}$ of Proposition~\ref{prop:warmup}. Its basic $25$-vertex layout is adapted from Fig.~1(a) of Svozil~\cite{svozil-2025-color}, preserving the threefold symmetry of the base; smooth dashed continuations connect the ends of its context traces to the common apex $26$, shown at the side. Figure~\ref{fig:h65-hypergraph} presents the full $65$-vertex, $29$-context four-uniform hypergraph diagram of the canonical linear hypergraph $H_{65}$.

% --- Styles and Colors from thicker-band-styles.tex ---
\definecolor{ucA01}{HTML}{1C638E}
\definecolor{ucA02}{HTML}{257D72}
\definecolor{ucA03}{HTML}{BC681D}
\definecolor{ucA04}{HTML}{8B3D88}
\definecolor{ucA05}{HTML}{82623E}
\definecolor{ucA06}{HTML}{0000FF}
\definecolor{ucA07}{HTML}{00B400}
\definecolor{ucA08}{HTML}{FF15F4}
\definecolor{ucA09}{HTML}{FF005F}
\definecolor{ucA10}{HTML}{006AFF}
\definecolor{ucA11}{HTML}{8A9400}
\definecolor{ucA12}{HTML}{FF3500}
\definecolor{ucA13}{HTML}{E97F94}
\definecolor{ucA14}{HTML}{005F0A}
\definecolor{ucA15}{HTML}{FF4AB4}
\definecolor{ucA16}{HTML}{740AB4}
\definecolor{ucA17}{HTML}{6A9FFF}
\definecolor{ucA18}{HTML}{941F2A}
\definecolor{ucA19}{HTML}{C97FFF}
\definecolor{ucA20}{HTML}{00B474}
\definecolor{ucA21}{HTML}{7F9F5F}
\definecolor{ucA22}{HTML}{B494C9}
\definecolor{ucA23}{HTML}{2A4A9F}
\definecolor{ucA24}{HTML}{FF6A55}
\definecolor{ucA25}{HTML}{AA005F}
\definecolor{ucA26}{HTML}{AA15FF}
\definecolor{ucA27}{HTML}{35AAC9}
\definecolor{ucA28}{HTML}{7F3F55}
\definecolor{ucA29}{HTML}{0035E9}
\definecolor{sgink}{HTML}{3D454B}
\definecolor{sggray}{HTML}{858B90}
\tikzset{
 sg/context/.style={line width=1.15pt,line cap=round,draw=sggray},
 sg/base/.style={line width=1.7pt,line cap=round},
 sg/label/.style={font=\sffamily\small,text=sgink,fill=white,inner sep=1pt},
 sg/number/.style={inner sep=0pt,outer sep=0pt,text=black,
   font=\sffamily\fontsize{8}{8}\selectfont},
}
\newcommand{\FCvertex}[2]{%
  \foreach \diskcolor/\diskradius in {#2}{%
    \fill[\diskcolor] (v#1) circle[radius=\diskradius cm];%
  }%
  \fill[white] (v#1) circle[radius=.160cm];%
  \node[sg/number] at (v#1) {#1};%
}

\begin{figure*}[t]
\centering
\resizebox{.90\linewidth}{!}{%
\begin{tikzpicture}[x=1cm,y=1cm]
% The nonagon is the symmetric outline of Fig. 1(a) of the earlier paper.
% Rotating by 120 degrees cycles the three coordinate directions.
\def\YOR{5.4}
\foreach \vertex/\angle in {1/70,3/110,11/150,7/190,9/230,10/270,4/310,6/350,13/30}{%
  \coordinate (v\vertex) at (\angle:\YOR);
}
\coordinate (v12) at (0,0); % h_0: fixed by the threefold rotation
% Three coordinate rays, exactly at the midpoints of the coordinate triads.
\coordinate (v2) at ($(v1)!.5!(v3)$);
\coordinate (v8) at ($(v7)!.5!(v9)$);
\coordinate (v5) at ($(v4)!.5!(v6)$);
% Six completion vertices on the outer nonagon.
\foreach \vertex/\first/\last in {14/3/11,15/11/7,23/9/10,18/10/4,22/6/13,20/13/1}{%
  \coordinate (v\vertex) at ($(v\first)!.5!(v\last)$);
}
% Three completions on the h_0 spokes.
\foreach \vertex/\last in {25/3,16/9,17/6}{%
  \coordinate (v\vertex) at ($(v12)!.52!(v\last)$);
}
% Three completions on the long chords: identical fractional positions.
\foreach \vertex/\first/\last in {24/11/4,19/10/1,21/13/7}{%
  \coordinate (v\vertex) at ($(v\first)!.75!(v\last)$);
}

% Dashed paths continue outward from terminal vertices, with exact tangency.
\coordinate (v26) at (8.05,3.55);
\coordinate (YOgrayBefore) at ($(v2)+(-1,0)$);
\YOapexlink{1}{v2}{v1}{0.351939419}{1.326540885}{(5.550000,4.250000)}
\YOapexlink{2}{v8}{v7}{0.351939419}{1.813841619}{(4.050000,6.550000)}
\YOapexlink{3}{v5}{v6}{0.351939419}{1.326540885}{(6.650000,3.550000)}
\YOapexlink{4}{YOgrayBefore}{v2}{0.650000000}{2.450000000}{(2.050000,2.050000)}
\YOapexlink{5}{v25}{v12}{0.231481481}{2.368233618}{(5.550000,6.550000)}
\YOapexlink{6}{v16}{v12}{0.231481481}{0.872507123}{(6.650000,2.050000)}
\YOapexlink{7}{v17}{v12}{0.231481481}{1.655982906}{(0.050000,6.550000)}
\YOapexlink{8}{v14}{v3}{0.351939419}{1.326540885}{(2.050000,6.550000)}
\YOapexlink{9}{v15}{v11}{0.351939419}{1.813841619}{(5.550000,2.050000)}
\YOapexlink{10}{v24}{v11}{0.081484852}{0.582930092}{(0.050000,8.550000)}
\YOapexlink{11}{v23}{v10}{0.351939419}{2.517720456}{(2.050000,2.050000)}
\YOapexvia{12}{v19}{v10}{0.081484852}{0.231918424}{(0.100000,-6.900000)}{(1.300000,-6.900000)}{(3.300000,-6.900000)}{(2.050000,2.050000)}
\YOapexlink{13}{v18}{v4}{0.351939419}{3.600610974}{(4.050000,2.050000)}
\YOapexlink{14}{v20}{v13}{0.351939419}{1.326540885}{(6.650000,2.850000)}
\YOapexlink{15}{v21}{v13}{0.081484852}{0.307135210}{(6.650000,3.550000)}
\YOapexlink{16}{v22}{v13}{0.351939419}{1.326540885}{(6.650000,4.250000)}

% Context 4: one smooth, closed threefold curve through precisely v2,v8,v5.
% Its three Bezier pieces are exact rotations of the same seed piece.
\pgfmathsetmacro{\YObasisR}{\YOR*cos(20)}
\foreach \rotation in {0,120,240}{%
  \begin{scope}[rotate=\rotation]
    \draw[YOtrace,YOctx4]
      (0,\YObasisR)
      .. controls ({-.5*\YObasisR},\YObasisR)
      and ({-(sqrt(3)/2+.25)*\YObasisR},{(-.5+.25*sqrt(3))*\YObasisR})
      .. ({-sqrt(3)/2*\YObasisR},{-.5*\YObasisR});
  \end{scope}
}
% Each solid trace is completed by its same-colour dashed connection to 26.
\YOline{1}{3}{2}{1}
\YOline{2}{9}{8}{7}
\YOline{3}{4}{5}{6}
\YOline{5}{3}{25}{12}
\YOline{6}{9}{16}{12}
\YOline{7}{6}{17}{12}
\YOline{8}{11}{14}{3}
\YOline{9}{7}{15}{11}
\YOline{10}{4}{24}{11}
\YOline{11}{9}{23}{10}
\YOline{12}{1}{19}{10}
\YOline{13}{10}{18}{4}
\YOline{14}{1}{20}{13}
\YOline{15}{7}{21}{13}
\YOline{16}{6}{22}{13}

% Nested context-colour rings, drawn after all traces.
\YOvertex{1}{1/.34,12/.29,14/.24}
\YOvertex{2}{1/.29,4/.24}
\YOvertex{3}{1/.34,5/.29,8/.24}
\YOvertex{4}{3/.34,10/.29,13/.24}
\YOvertex{5}{3/.29,4/.24}
\YOvertex{6}{3/.34,7/.29,16/.24}
\YOvertex{7}{2/.34,15/.29,9/.24}
\YOvertex{8}{2/.29,4/.24}
\YOvertex{9}{2/.34,6/.29,11/.24}
\YOvertex{10}{11/.34,13/.29,12/.24}
\YOvertex{11}{8/.34,9/.29,10/.24}
\YOvertex{12}{5/.34,6/.29,7/.24}
\YOvertex{13}{16/.34,14/.29,15/.24}
\YOvertex{14}{8/.24}
\YOvertex{15}{9/.24}
\YOvertex{16}{6/.24}
\YOvertex{17}{7/.24}
\YOvertex{18}{13/.24}
\YOvertex{19}{12/.24}
\YOvertex{20}{14/.24}
\YOvertex{21}{15/.24}
\YOvertex{22}{16/.24}
\YOvertex{23}{11/.24}
\YOvertex{24}{10/.24}
\YOvertex{25}{5/.24}

% Traditional symbols are supplementary labels; all numbers match the main paper.
\foreach \angle/\symbol in {70/{y_1^+},110/{y_1^-},190/{y_2^+},230/{y_2^-},310/{y_3^+},350/{y_3^-}}{%
  \node[YOlabel] at (\angle:6.02) {$\symbol$};
}
\node[YOlabel,anchor=east] at ($(v11)+(-.45,-.22)$) {$h_1$};
\node[YOlabel] at ($(v10)+(-.48,.60)$) {$h_2$};
\node[YOlabel] at ($(v13)+(-.76,.12)$) {$h_3$};
\foreach \angle/\symbol in {90/{z_1},210/{z_2},330/{z_3}}{%
  \node[YOlabel] at (\angle:{\YObasisR-.58}) {$\symbol$};
}
\node[YOlabel] at (.16,.62) {$h_0$};

% Common apex: dashed context continuations meet underneath these rings.
\draw[draw=black!20,rounded corners=5pt] (6.35,1.50) rectangle (9.75,5.10);
\node[font=\sffamily\small] at (8.05,4.72) {Four-dimensional lift};
\foreach \context/\radius in {16/.88,15/.84,14/.80,13/.76,12/.72,11/.68,10/.64,9/.60,8/.56,7/.52,6/.48,5/.44,4/.40,3/.36,2/.32,1/.28}{%
  \fill[YOctx\context] (v26) circle[radius=\radius cm];
}
\fill[white] (v26) circle[radius=.22cm];
\node[YOnumber] at (v26) {26};
\node[font=\small] at (8.05,2.28) {common apex};
\node[font=\small] at (8.05,1.80) {$e_m^{(4)}=e_m\cup\{26\}$};
% Use the visible curve extents, excluding distant Bezier control handles.
\pgfresetboundingbox
\path[use as bounding box] (-6.580000,-7.120032) rectangle (10.080000,6.480000);
\end{tikzpicture}%
}
\caption{Four-uniform context diagram of the $26$-vertex, $16$-context apex suspension $H_{\YO}^{(4)}$ of Proposition~\ref{prop:warmup}. The basic $25$-vertex Yu--Oh layout is redrawn from Svozil~\cite[Fig.~1(a)]{svozil-2025-color}, based on the configuration of Yu and Oh~\cite{Yu-2012}; the vertex numbering is that of this paper. The base geometry and context colors retain the $120^\circ$ rotational symmetry. Each solid trace contains one three-element base context; its dashed continuation leaves a terminal vertex in the same straight direction and bends smoothly to the common apex $26$, completing the four-element context $e_m^{(4)}=e_m\cup\{26\}$. The gray coordinate context $\{2,5,8\}$ continues tangentially at vertex $2$. Colors distinguish six context orbits under the rotation, and concentric rings record incident context colors. The supplementary labels $z_i,y_i^\pm$ ($i=1,2,3$) and $h_j$ ($j=0,1,2,3$) are the conventional names of the original $13$ Yu--Oh rays. Unmarked crossings are not vertices. Coordinates and all $16$ contexts are given in Eq.~\eqref{eq:hyo4-vectors} and Table~\ref{tab:app-hyo4-contexts}.}\label{fig:hyo4}
\end{figure*}

\begin{figure*}[tbp]
\centering
\resizebox{.90\linewidth}{!}{%
\begin{tikzpicture}[x=1cm,y=1cm]
\coordinate (v1) at (-2.000000,-2.000000);
\coordinate (v2) at (-5.000000,5.000000);
\coordinate (v3) at (5.000000,-5.000000);
\coordinate (v4) at (2.000000,2.000000);
\coordinate (v5) at (-5.000000,-5.000000);
\coordinate (v6) at (-2.000000,2.000000);
\coordinate (v7) at (2.000000,-2.000000);
\coordinate (v8) at (5.000000,5.000000);
\coordinate (v9) at (-1.521690,0.927051);
\coordinate (v10) at (1.800000,-0.860561);
\coordinate (v11) at (-1.521690,-0.927051);
\coordinate (v12) at (0.000000,7.300000);
\coordinate (v13) at (0.850000,0.000000);
\coordinate (v14) at (0.000000,2.333000);
\coordinate (v15) at (0.000000,-2.333000);
\coordinate (v16) at (2.333000,0.000000);
\coordinate (v17) at (-2.333000,0.000000);
\coordinate (v18) at (-1.700000,6.518000);
\coordinate (v19) at (-3.350000,5.759000);
\coordinate (v20) at (7.510310,2.131747);
\coordinate (v21) at (7.443563,-1.627935);
\coordinate (v22) at (0.329875,-0.365000);
\coordinate (v23) at (-0.204500,-0.740000);
\coordinate (v24) at (1.172000,0.560000);
\coordinate (v25) at (1.499750,1.130000);
\coordinate (v26) at (-2.225000,3.519815);
\coordinate (v27) at (-3.575000,4.239905);
\coordinate (v28) at (0.665000,2.222278);
\coordinate (v29) at (1.325000,2.112388);
\coordinate (v30) at (-0.800000,-2.199800);
\coordinate (v31) at (-1.400000,-2.099900);
\coordinate (v32) at (2.225000,-3.519815);
\coordinate (v33) at (3.612500,-4.259908);
\coordinate (v34) at (3.579823,-2.337500);
\coordinate (v35) at (4.099888,-3.312500);
\coordinate (v36) at (2.222278,0.665000);
\coordinate (v37) at (2.112388,1.325000);
\coordinate (v38) at (-2.189810,-0.860000);
\coordinate (v39) at (-2.094905,-1.430000);
\coordinate (v40) at (-3.299787,1.812500);
\coordinate (v41) at (-3.919865,2.975000);
\coordinate (v42) at (-7.534215,1.984465);
\coordinate (v43) at (-7.415219,-1.760849);
\coordinate (v44) at (1.400000,6.656000);
\coordinate (v45) at (2.787500,6.017750);
\coordinate (v46) at (0.116125,0.515000);
\coordinate (v47) at (-0.632000,1.040000);
\coordinate (v48) at (1.120250,-0.470000);
\coordinate (v49) at (1.646375,-1.385000);
\coordinate (v50) at (-0.515000,2.247253);
\coordinate (v51) at (-1.040000,2.159840);
\coordinate (v52) at (2.225000,3.519815);
\coordinate (v53) at (3.612500,4.259908);
\coordinate (v54) at (-2.225000,-3.519815);
\coordinate (v55) at (-3.612500,-4.259908);
\coordinate (v56) at (0.665000,-2.222278);
\coordinate (v57) at (1.325000,-2.112388);
\coordinate (v58) at (2.149850,-1.100000);
\coordinate (v59) at (2.076590,-1.540000);
\coordinate (v60) at (3.319790,1.850000);
\coordinate (v61) at (3.939868,3.012500);
\coordinate (v62) at (-3.319790,-1.850000);
\coordinate (v63) at (-3.939868,-3.012500);
\coordinate (v64) at (-2.222278,0.665000);
\coordinate (v65) at (-2.112388,1.325000);
\draw[sg/context,draw=ucA06] (v12) -- (v2); % E1
\draw[sg/context,draw=ucA07] (0.000000,7.300000) .. controls (10.000000,9.000000) and (8.500000,-5.000000) .. (5.000000,-5.000000); % E2
\draw[sg/context,draw=ucA08] (v13) -- (v1); % E3
\draw[sg/context,draw=ucA09] (v13) -- (v4); % E4
\draw[sg/context,draw=ucA10] (v14) -- (v2); % E5
\draw[sg/context,draw=ucA11] (v14) -- (v4); % E6
\draw[sg/context,draw=ucA12] (v15) -- (v1); % E7
\draw[sg/context,draw=ucA13] (v15) -- (v3); % E8
\draw[sg/context,draw=ucA14] (v16) -- (v3); % E9
\draw[sg/context,draw=ucA15] (v16) -- (v4); % E10
\draw[sg/context,draw=ucA16] (v17) -- (v1); % E11
\draw[sg/context,draw=ucA17] (v17) -- (v2); % E12
\draw[sg/context,draw=ucA18] (0.000000,7.300000) .. controls (-10.000000,9.000000) and (-8.500000,-5.000000) .. (-5.000000,-5.000000); % E13
\draw[sg/context,draw=ucA19] (v12) -- (v8); % E14
\draw[sg/context,draw=ucA20] (v13) -- (v6); % E15
\draw[sg/context,draw=ucA21] (v13) -- (v7); % E16
\draw[sg/context,draw=ucA22] (v14) -- (v6); % E17
\draw[sg/context,draw=ucA23] (v14) -- (v8); % E18
\draw[sg/context,draw=ucA24] (v15) -- (v5); % E19
\draw[sg/context,draw=ucA25] (v15) -- (v7); % E20
\draw[sg/context,draw=ucA26] (v16) -- (v7); % E21
\draw[sg/context,draw=ucA27] (v16) -- (v8); % E22
\draw[sg/context,draw=ucA28] (v17) -- (v5); % E23
\draw[sg/context,draw=ucA29] (v17) -- (v6); % E24
\draw[sg/base,draw=ucA01,preaction={draw=white,line width=3.0pt}] (-5.000000,5.000000) .. controls (-6.104569,3.895431) and (-4.761424,0.761424) .. (-2.000000,-2.000000) .. controls (0.761424,-4.761424) and (3.895431,-6.104569) .. (5.000000,-5.000000) .. controls (6.104569,-3.895431) and (4.761424,-0.761424) .. (2.000000,2.000000); % C1
\draw[sg/base,draw=ucA02,preaction={draw=white,line width=3.0pt}] (-5.000000,-5.000000) .. controls (-6.104569,-3.895431) and (-4.761424,-0.761424) .. (-2.000000,2.000000) .. controls (0.761424,4.761424) and (3.895431,6.104569) .. (5.000000,5.000000) .. controls (6.104569,3.895431) and (4.761424,0.761424) .. (2.000000,-2.000000); % C2
\draw[sg/base,draw=ucA03,preaction={draw=white,line width=3.0pt}] (0.000000,7.300000) .. controls (-0.507230,5.175684) and (-1.521690,1.593718) .. (-1.521690,0.927051) .. controls (-1.521690,0.260384) and (0.183333,0.000000) .. (0.850000,0.000000) .. controls (1.516667,0.000000) and (1.133333,-0.427228) .. (1.800000,-0.860561); % C3
\draw[sg/base,draw=ucA04,preaction={draw=white,line width=3.0pt}] (0.000000,2.333000) .. controls (0.000000,1.883000) and (-1.307492,2.163119) .. (-1.521690,0.927051) .. controls (-1.626103,0.324521) and (-1.626103,-0.324521) .. (-1.521690,-0.927051) .. controls (-1.307492,-2.163119) and (0.000000,-1.883000) .. (-0.000000,-2.333000); % C4
\draw[sg/base,draw=ucA05,preaction={draw=white,line width=3.0pt}] (2.333000,0.000000) .. controls (1.733000,0.000000) and (2.555418,-0.657410) .. (1.800000,-0.860561) .. controls (0.831586,-1.120992) and (-0.478454,-1.147215) .. (-1.521690,-0.927051) .. controls (-2.437650,-0.733747) and (-1.733000,0.000000) .. (-2.333000,-0.000000); % C5
\FCvertex{1}{ucA16/0.348000,ucA12/0.302000,ucA08/0.256000,ucA01/0.210000}
\FCvertex{2}{ucA17/0.348000,ucA10/0.302000,ucA06/0.256000,ucA01/0.210000}
\FCvertex{3}{ucA14/0.348000,ucA13/0.302000,ucA07/0.256000,ucA01/0.210000}
\FCvertex{4}{ucA15/0.348000,ucA11/0.302000,ucA09/0.256000,ucA01/0.210000}
\FCvertex{5}{ucA28/0.348000,ucA24/0.302000,ucA18/0.256000,ucA02/0.210000}
\FCvertex{6}{ucA29/0.348000,ucA22/0.302000,ucA20/0.256000,ucA02/0.210000}
\FCvertex{7}{ucA26/0.266111,ucA25/0.240148,ucA21/0.214185,ucA02/0.188221}
\FCvertex{8}{ucA27/0.348000,ucA23/0.302000,ucA19/0.256000,ucA02/0.210000}
\FCvertex{9}{ucA04/0.256000,ucA03/0.210000}
\FCvertex{10}{ucA05/0.220454,ucA03/0.191487}
\FCvertex{11}{ucA05/0.256000,ucA04/0.210000}
\FCvertex{12}{ucA19/0.394000,ucA18/0.348000,ucA07/0.302000,ucA06/0.256000,ucA03/0.210000}
\FCvertex{13}{ucA21/0.333158,ucA20/0.299119,ucA09/0.265079,ucA08/0.231039,ucA03/0.197000}
\FCvertex{14}{ucA23/0.316623,ucA22/0.285834,ucA11/0.255045,ucA10/0.224256,ucA04/0.193466}
\FCvertex{15}{ucA25/0.394000,ucA24/0.348000,ucA13/0.302000,ucA12/0.256000,ucA04/0.210000}
\FCvertex{16}{ucA27/0.394000,ucA26/0.348000,ucA15/0.302000,ucA14/0.256000,ucA05/0.210000}
\FCvertex{17}{ucA29/0.394000,ucA28/0.348000,ucA17/0.302000,ucA16/0.256000,ucA05/0.210000}
\FCvertex{18}{ucA06/0.210000}
\FCvertex{19}{ucA06/0.210000}
\FCvertex{20}{ucA07/0.210000}
\FCvertex{21}{ucA07/0.210000}
\FCvertex{22}{ucA08/0.210000}
\FCvertex{23}{ucA08/0.210000}
\FCvertex{24}{ucA09/0.210000}
\FCvertex{25}{ucA09/0.210000}
\FCvertex{26}{ucA10/0.210000}
\FCvertex{27}{ucA10/0.210000}
\FCvertex{28}{ucA11/0.210000}
\FCvertex{29}{ucA11/0.210000}
\FCvertex{30}{ucA12/0.210000}
\FCvertex{31}{ucA12/0.210000}
\FCvertex{32}{ucA13/0.210000}
\FCvertex{33}{ucA13/0.210000}
\FCvertex{34}{ucA14/0.210000}
\FCvertex{35}{ucA14/0.210000}
\FCvertex{36}{ucA15/0.210000}
\FCvertex{37}{ucA15/0.210000}
\FCvertex{38}{ucA16/0.210000}
\FCvertex{39}{ucA16/0.210000}
\FCvertex{40}{ucA17/0.210000}
\FCvertex{41}{ucA17/0.210000}
\FCvertex{42}{ucA18/0.210000}
\FCvertex{43}{ucA18/0.210000}
\FCvertex{44}{ucA19/0.210000}
\FCvertex{45}{ucA19/0.210000}
\FCvertex{46}{ucA20/0.210000}
\FCvertex{47}{ucA20/0.210000}
\FCvertex{48}{ucA21/0.197000}
\FCvertex{49}{ucA21/0.210000}
\FCvertex{50}{ucA22/0.193466}
\FCvertex{51}{ucA22/0.210000}
\FCvertex{52}{ucA23/0.210000}
\FCvertex{53}{ucA23/0.210000}
\FCvertex{54}{ucA24/0.210000}
\FCvertex{55}{ucA24/0.210000}
\FCvertex{56}{ucA25/0.210000}
\FCvertex{57}{ucA25/0.210000}
\FCvertex{58}{ucA26/0.191487}
\FCvertex{59}{ucA26/0.188221}
\FCvertex{60}{ucA27/0.210000}
\FCvertex{61}{ucA27/0.210000}
\FCvertex{62}{ucA28/0.210000}
\FCvertex{63}{ucA28/0.210000}
\FCvertex{64}{ucA29/0.210000}
\FCvertex{65}{ucA29/0.210000}
\node[sg/label,text=ucA01,anchor=south east] at (-4.70,5.65) {$C_1$};
\node[sg/label,text=ucA02,anchor=north east] at (-4.70,-5.65) {$C_2$};
\node[sg/label,text=ucA03,anchor=south] at (0,7.86) {$C_3$};
\node[sg/label,text=ucA04,anchor=south] at (0,2.90) {$C_4$};
\node[sg/label,text=ucA05,anchor=west] at (2.94,0) {$C_5$};
\end{tikzpicture}%
}
\caption{Complete four-uniform hypergraph diagram of the canonical linear $4$-uniform hypergraph $H_{65}=(V_{65},E_{29})$ in $\R^4$ of Theorem~\ref{thm:h65} ($65$ pairwise non-collinear integer vectors $v_1,\ldots,v_{65}\in\mathbb{Z}^4$ enumerated in Eq.~\eqref{eq:h65-coordinates} and Table~\ref{tab:app-h65-vectors}, $29$ four-element contexts listed in Eq.~\eqref{eq:h65-contexts} and Table~\ref{tab:app-h65-contexts}, and $29\times\binom{4}{2}=174$ orthogonality edges). All $29$ four-element contexts have distinct colors, matched by the concentric filled disk layers of successive radii surrounding each vertex number, with a white center behind every label: the five core contexts $C_1=\{1,2,3,4\}$, $C_2=\{5,6,7,8\}$, $C_3=\{9,10,12,13\}$, $C_4=\{9,11,14,15\}$, and $C_5=\{10,11,16,17\}$ are drawn as bold smooth curves, while the $24$ bridge-completion contexts $E_1,\ldots,E_{24}$ (such as the straight line $E_{21}=\{7,16,58,59\}$) connect the binary bridge vertices $12,\ldots,17$ (each surrounded by $5$ nested colored disks) to the basis vertices $1,\ldots,8$ ($4$ nested disks) via the degree-$1$ completion pairs $18,\ldots,65$ ($1$ colored disk), with the hub vertices $9,10,11$ surrounded by $2$ nested disks. Unmarked crossings are not vertices.}\label{fig:h65-hypergraph}
\end{figure*}
Both drawings represent $4$-uniform context hypergraphs. They differ from Greechie diagrams in the strict sense of Ref.~\cite[Definition~2.4]{svozil-tkadlec}: the warm-up has contexts sharing two vertices, whereas $H_{65}$ contains the three-context Berge cycle $C_3,C_4,C_5$, with successive intersection vertices $9,11,10$. Strict Greechie diagrams permit neither feature. The smooth-curve drawing convention is used only to display the listed incidences.


\begin{table*}[t]
\caption{The $24$ core state supports $S_\alpha\subset V_{\mathrm{core}}$ ($\alpha=1,\ldots,24$) of $H_{65}$, grouped by their hub ray $h\in\{9,10,11\}$ and opposite bridge ray $b\in B_1\cup B_2\cup B_3$.}\label{tab:core-supports}
\begin{ruledtabular}
\begin{tabular}{cccc}
Hub $h$ & Bridge $b$ & Allowed $(i_1,i_2)\in C_1\times C_2$ & Supports $S_\alpha=\{i_1,i_2,h,b\}$ \\ \hline
$9$  & $16$ & $\{1,2\}\times\{5,6\}$ & $S_1=\{1,5,9,16\}$,\ \ $S_3=\{1,6,9,16\}$,\ \ $S_7=\{2,5,9,16\}$,\ \ $S_9=\{2,6,9,16\}$ \\
$9$  & $17$ & $\{3,4\}\times\{7,8\}$ & $S_{15}=\{3,7,9,17\}$,\ \ $S_{17}=\{3,8,9,17\}$,\ \ $S_{21}=\{4,7,9,17\}$,\ \ $S_{23}=\{4,8,9,17\}$ \\ \hline
$10$ & $14$ & $\{1,3\}\times\{5,7\}$ & $S_2=\{1,5,10,14\}$,\ \ $S_5=\{1,7,10,14\}$,\ \ $S_{13}=\{3,5,10,14\}$,\ \ $S_{16}=\{3,7,10,14\}$ \\
$10$ & $15$ & $\{2,4\}\times\{6,8\}$ & $S_{10}=\{2,6,10,15\}$,\ \ $S_{11}=\{2,8,10,15\}$,\ \ $S_{19}=\{4,6,10,15\}$,\ \ $S_{24}=\{4,8,10,15\}$ \\ \hline
$11$ & $12$ & $\{1,4\}\times\{6,7\}$ & $S_4=\{1,6,11,12\}$,\ \ $S_6=\{1,7,11,12\}$,\ \ $S_{20}=\{4,6,11,12\}$,\ \ $S_{22}=\{4,7,11,12\}$ \\
$11$ & $13$ & $\{2,3\}\times\{5,8\}$ & $S_8=\{2,5,11,13\}$,\ \ $S_{12}=\{2,8,11,13\}$,\ \ $S_{14}=\{3,5,11,13\}$,\ \ $S_{18}=\{3,8,11,13\}$
\end{tabular}
\end{ruledtabular}
\end{table*}

\section{Main theorem and proofs}

\begin{theorem}[Properties of the canonical linear $4$-uniform hypergraph $H_{65}$]\label{thm:h65}
The $65$-vertex, $29$-context $4$-uniform hypergraph $H_{65}=(V_{65},E_{29})$ defined by Eqs.~\eqref{eq:h65-coordinates}--\eqref{eq:h65-contexts} satisfies:
\begin{enumerate}
\item \emph{Linearity and clique number:} $H_{65}$ is linear ($|e\cap e'|\le1$ for all distinct $e,e'\in E_{29}$), and the $4$-cliques of $\Gamma(H_{65})$ are precisely the $29$ contexts of $E_{29}$; in particular, $\omega(\Gamma(H_{65}))=4$. It is nonconformal, with $13$ additional maximal triangles.
\item \emph{Edge-faithful realization in $\R^4$:} The $65$ integer vectors $v_1,\ldots,v_{65}\in\R^4$ in Eq.~\eqref{eq:h65-coordinates} are pairwise non-collinear and satisfy $v_i\cdot v_j=0$ if and only if $\{i,j\}\in E(\Gamma(H_{65}))$ ($174$ orthogonal pairs).
\item \emph{Chromatic--clique gap $\chi(H_{65})=5>4=\omega(\Gamma(H_{65}))$:} The $17$-vertex core $G_{\mathrm{core}}$ has independence number $\alpha(G_{\mathrm{core}})=4$, fractional chromatic number $\chi_f(\Gamma(H_{65}))=17/4=4.25>4$, and $\chi(H_{65})=5$.
\item \emph{Set-representability:} The core $(V_{\mathrm{core}},\{C_1,\ldots,C_5\})$ admits exactly $24$ two-valued state supports $S_1,\ldots,S_{24}$ compatible with $G_{\mathrm{core}}$ (listed in Table~\ref{tab:core-supports}), and the resulting $48$ canonical two-valued states $s_{\alpha,\varepsilon}$ on $H_{65}$ via Eq.~\eqref{eq:canonical-states} form a unital and separating collection on all $65$ vertices (within a full state space of $24\times 2^{14}=393{,}216$ states, enumerated in Appendix~\ref{app:states-h65}).
\end{enumerate}
\end{theorem}

\begin{proof}
(1)--(2) By inspection of Eq.~\eqref{eq:h65-contexts}, the five base contexts $C_1,\ldots,C_5$ satisfy $|C_a\cap C_b|\le1$ ($C_1,C_2$ are disjoint from everything else, while $C_3\cap C_4=\{9\}$, $C_3\cap C_5=\{10\}$, $C_4\cap C_5=\{11\}$), and the $24$ pairs $e_m=\{u_m,v_m\}$ in $E_1,\ldots,E_{24}$ are precisely the $24$ orthogonal pairs between $C_1\cup C_2$ and $B_1\cup B_2\cup B_3$.

To determine the triangles and $4$-cliques of $G_{\mathrm{core}}$ and $\Gamma(H_{65})$, let $e_m=\{u_m,v_m\}$ be a bridge edge with $u_m\in C_a$ ($a\in\{1,2\}$) and $v_m\in B_k$. Because $u_m$ is non-orthogonal to the other basis $C_{3-a}$ and to $\{9,10,11\}$, while $v_m$ has no neighbors in $(B_1\cup B_2\cup B_3)\setminus B_k$ and its partner $v'_m\in B_k$ is non-orthogonal to $u_m$ (Table~\ref{tab:bridge-neighbors}), the \emph{unique} common third neighbor of $u_m$ and $v_m$ in $G_{\mathrm{core}}$ is the second vertex of $N(v_m)\cap C_a$. Consequently, $G_{\mathrm{core}}$ has exactly $13$ triangles not contained in $C_1,\ldots,C_5$: the $12$ bridge triangles
\begin{align*}
&\{1,2,17\},\ \{1,3,15\},\ \{1,4,13\},\\
&\{2,3,12\},\ \{2,4,14\},\ \{3,4,16\},\\
&\{5,6,17\},\ \{5,7,15\},\ \{5,8,12\},\\
&\{6,7,13\},\ \{6,8,14\},\ \{7,8,16\},
\end{align*}
together with the hub triad $\{9,10,11\}$. Since each bridge edge has only one common third neighbor in $V_{\mathrm{core}}$, \emph{no $4$-clique ($K_4$) in $G_{\mathrm{core}}$ contains a bridge edge}. This has two immediate consequences:
\begin{itemize}
\item First, consider two distinct bridge edges $e_\ell,e_m$. If they have four distinct endpoints, orthogonality $L_\ell\perp L_m$ would make those endpoints a $K_4$ containing a bridge edge, which was just excluded. If they share an endpoint $x$, orthogonality of the spans would imply $v_x\cdot v_x=0$, impossible for a nonzero real vector. Thus $L_\ell\not\perp L_m$ ($P_\ell\ne L_m$), verifying Eq.~\eqref{eq:plane-nonorth} in both cases.

To verify Eq.~\eqref{eq:span-condition}, consider $a v_{u_m}+b v_{v_m}$ with $ab\ne0$. If $u_m\in C_1$, its nonzero coordinate lies outside the two-coordinate support of $v_{v_m}$, so the combination has support size three; no core vector has that support size. If $u_m\in C_2$, write $v_{u_m}=h\in\{\pm1\}^4$. Orthogonality makes the two nonzero components of $v_{v_m}$ have opposite signs relative to those of $h$. After rescaling, the component magnitudes are those of
\begin{equation*}
(1+t,\,1-t,\,1,\,1),\qquad t\in\R\setminus\{0\},
\end{equation*}
up to permutation. For $t=\pm1$ the support size is three. Otherwise all four components are nonzero, so proportionality to a core vector would require $|1+t|=|1-t|=1$, which forces $t=0$. Both cases exclude any other core ray from $L_m$. The reproducible completion rule and exact integer verification in Appendix~\ref{app:verification} confirm that $v_i\cdot v_j=0$ holds exactly for the $5\times 6+24\times 6=174$ edges of $\Gamma(H_{65})$.
\item Second, deleting all $24$ bridge edges from $G_{\mathrm{core}}$ leaves only the edges of $C_1,\ldots,C_5$, whose only $K_4$ cliques are $C_1,\ldots,C_5$ themselves. Since every fresh completion vertex $w\in\{18,\ldots,65\}$ has degree $3$ in $\Gamma(H_{65})$ with neighborhood equal to the remaining three vertices of its unique completion context $E_m$, the $4$-cliques of $\Gamma(H_{65})$ are precisely the $29$ contexts of $E_{29}$, and $\omega(\Gamma(H_{65}))=4$ (which also follows directly from the injective orthogonal realization in $\R^4$).
\end{itemize}

\begin{table}[t]
\caption{Orthogonal neighbors in $C_1=\{1,2,3,4\}$ and $C_2=\{5,6,7,8\}$ for each of the six binary bridge rays $b\in B_1\cup B_2\cup B_3$.}\label{tab:bridge-neighbors}
\begin{ruledtabular}
\begin{tabular}{cccc}
Pair & Ray $b$ & $N(b)\cap C_1$ & $N(b)\cap C_2$ \\ \hline
$B_1$ & $12=(1,0,0,-1)$ & $\{2,3\}$ & $\{5,8\}$ \\
$B_1$ & $13=(0,1,1,0)$  & $\{1,4\}$ & $\{6,7\}$ \\ \hline
$B_2$ & $14=(1,0,1,0)$  & $\{2,4\}$ & $\{6,8\}$ \\
$B_2$ & $15=(0,1,0,-1)$ & $\{1,3\}$ & $\{5,7\}$ \\ \hline
$B_3$ & $16=(1,1,0,0)$  & $\{3,4\}$ & $\{7,8\}$ \\
$B_3$ & $17=(0,0,1,-1)$ & $\{1,2\}$ & $\{5,6\}$
\end{tabular}
\end{ruledtabular}
\end{table}

(3) To prove $\alpha(G_{\mathrm{core}})=4$ by hand, let $I\subseteq V_{\mathrm{core}}$ be an independent set in $G_{\mathrm{core}}$, and write $W_9=C_3\cup C_4\cup C_5=\{9,10,11\}\cup B_1\cup B_2\cup B_3$. Since $C_1$ and $C_2$ are $4$-cliques, $|I\cap C_1|\le1$ and $|I\cap C_2|\le1$. Moreover, $W_9$ is partitioned into the triangle $\{9,10,11\}$ and the three edges $B_1,B_2,B_3$, so $|I\cap W_9|\le 4$; in fact, if $I\cap\{9,10,11\}\ne\emptyset$, say $9\in I$, then $9\in C_3\cap C_4$ excludes all of $\{10,11,12,13,14,15\}$ and leaves at most one vertex in $B_3=\{16,17\}$, so $|I\cap W_9|\le2$ and hence $|I|\le 1+1+2=4$.

The only remaining way to have $|I|\ge 5$ is for $I\cap\{9,10,11\}=\emptyset$ and $|I\cap W_9|=3$, which requires $I\cap W_9=\{b_1,b_2,b_3\}$ with $(b_1,b_2,b_3)\in B_1\times B_2\times B_3$. Table~\ref{tab:bridge-neighbors} lists the orthogonal neighbors of each $b\in B_1\cup B_2\cup B_3$ in $C_1$ and $C_2$. Call a triple $(b_1,b_2,b_3)\in B_1\times B_2\times B_3$ even or odd according to the parity of the number of its entries in $\{13,15,17\}$, the second elements of the three ordered bridge pairs. Observe the complementary parity of the $2+2$ partitions of $C_1$ and $C_2$:
\begin{itemize}
\item For the four \emph{even-parity} triples $(b_1,b_2,b_3)\in\{(12,14,16),\allowbreak (12,15,17),\allowbreak (13,14,17),\allowbreak (13,15,16)\}$, the union of their orthogonal neighbors in $C_2$ is all of $C_2=\{5,6,7,8\}$ (leaving only one non-neighbor in $C_1$ and \emph{zero} non-neighbors in $C_2$), so $|I|\le 3+1+0=4$.
\item For the four \emph{odd-parity} triples $(b_1,b_2,b_3)\in\{(12,14,17),\allowbreak (12,15,16),\allowbreak (13,14,16),\allowbreak (13,15,17)\}$, the union of their orthogonal neighbors in $C_1$ is all of $C_1=\{1,2,3,4\}$ (leaving \emph{zero} non-neighbors in $C_1$ and only one non-neighbor in $C_2$), so again $|I|\le 3+0+1=4$.
\end{itemize}
Thus $\alpha(G_{\mathrm{core}})=4$ (attained, e.g., by $S_1=\{1,5,9,16\}$). In any proper $k$-coloring of $G_{\mathrm{core}}$, each color class is an independent set of size at most $4$, so $4k\ge |V_{\mathrm{core}}|=17$, forcing $k\ge\lceil 17/4\rceil=5$ and $\chi_f(\Gamma(H_{65}))\ge\chi_f(G_{\mathrm{core}})\ge 17/4=4.25$. Together with the explicit $5$-coloring in Eq.~\eqref{eq:h65-coordinates}, this establishes $\chi(H_{65})=5>4=\omega(\Gamma(H_{65}))$.

For the matching fractional upper bound, the independent four-sets of $G_{\mathrm{core}}$ consist of the $24$ sets in Table~\ref{tab:core-supports}, each with one hub, one bridge vertex, and one vertex from each of $C_1,C_2$; $12$ sets with two bridge vertices from different pairs and one vertex from each basis; and $8$ sets with three bridge vertices and one basis vertex. These counts follow from the neighbor lists in Table~\ref{tab:bridge-neighbors}, which also show that assigning weight $1/8$ to each set of the first type and $1/16$ to each of the other types covers every vertex with total weight one. The total weight is $24/8+12/16+8/16=17/4$.

Equivalently, take two copies of each of the first $24$ independent sets and one copy of each of the other $20$, obtaining $68$ color classes in which every core vertex occurs $16$ times. For any bridge edge, its adjacent endpoints occupy $32$ disjoint classes, leaving $36$ classes available. Assign $16$ available classes to its first completion vertex and $16$ other available classes to its second. Completion vertices from different contexts are nonadjacent, so these choices are independent. This extends the coloring to all $65$ vertices and proves $\chi_f(\Gamma(H_{65}))\le68/16=17/4$. Together with the lower bound, this proves equality.


(4) A subset $S\subset V_{\mathrm{core}}$ is the support of a core state compatible with $G_{\mathrm{core}}$ if and only if $S$ is independent in $G_{\mathrm{core}}$ and $|S\cap C_a|=1$ for $a=1,\ldots,5$. (Without the $24$ bridge-edge constraints, the five-context hypergraph $(V_{\mathrm{core}},\{C_1,\ldots,C_5\})$ alone has $4\times 4\times(3\times 2+2^3)=224$ two-valued states.) By the parity analysis in Part~(3), no triple $(b_1,b_2,b_3)\in B_1\times B_2\times B_3$ can be extended to include both a vertex of $C_1$ and a vertex of $C_2$. Consequently, $S$ must select:
\begin{itemize}
\item exactly one hub vertex $h\in\{9,10,11\}$ ($3$ choices), which satisfies the two contexts among $C_3,C_4,C_5$ containing $h$;
\item one binary bridge vertex $b$ from the unique remaining pair $B_k$ ($2$ choices), satisfying the third context among $C_3,C_4,C_5$;
\item one of the two vertices $i_1\in C_1\setminus N(b)$ ($2$ choices); and
\item one of the two vertices $i_2\in C_2\setminus N(b)$ ($2$ choices).
\end{itemize}
This yields precisely $3\times 2\times 2\times 2=24$ graph-compatible core state supports $S_1,\ldots,S_{24}$, listed in Table~\ref{tab:core-supports}.

To count the full state space $\mathcal{S}(H_{65})$, note that in the $24$-edge bipartite bridge graph between $C_1\cup C_2$ and $B_1\cup B_2\cup B_3$, every coordinate vertex $i_1\in C_1$ has bridge degree $3$, every Hadamard vertex $i_2\in C_2$ has bridge degree $3$, every bridge vertex $b\in B_1\cup B_2\cup B_3$ has bridge degree $4$, and every hub vertex $h\in\{9,10,11\}$ has bridge degree $0$. Because each $S_\alpha=\{i_1,i_2,h,b\}$ is independent in $G_{\mathrm{core}}$, no bridge edge has both endpoints in $S_\alpha$; hence \emph{every} $S_\alpha$ meets $3+3+0+4=10$ bridge edges and misses $24-10=14$ bridge edges. Each of the $14$ missed completion contexts $E_m=\{u_m,v_m,w_{m,1},w_{m,2}\}$ admits $2$ independent binary extensions $(s(w_{m,1}),s(w_{m,2}))\in\{(1,0),(0,1)\}$, while the $10$ hit completion contexts force $(0,0)$, giving $|\mathcal{S}(H_{65})|=24\times 2^{14}=393{,}216$.

Finally, to see that the $24$ edge zero-profiles $Z(e_m)$ in Eq.~\eqref{eq:zero-profiles} are nonempty and pairwise distinct, partition $S_1,\ldots,S_{24}$ into three blocks $\mathcal{B}_9,\mathcal{B}_{10},\mathcal{B}_{11}$ of $8$ supports each according to their hub ray $h\in\{9,10,11\}$ (the three horizontal sections of Table~\ref{tab:core-supports}). Every vertex $u\in C_1\cup C_2$ occurs in $2$ supports in each of the three blocks $\mathcal{B}_9,\mathcal{B}_{10},\mathcal{B}_{11}$ (total weight $6$, block weights $(2,2,2)$), whereas each bridge vertex $v\in B_k$ occurs in $4$ supports concentrated entirely in the single block $\mathcal{B}_{h_k}$ of the hub $h_k$ opposite $B_k$ (block weights $(4,0,0)$ up to permutation). Since $u_m\perp v_m$, the indicators $\mathbf{1}_{u_m}$ and $\mathbf{1}_{v_m}$ are disjoint, so the complement of $Z(e_m)$ has indicator $f_m=\mathbf{1}_{u_m}+\mathbf{1}_{v_m}$ of weight $6+4=10$ (so $|Z(e_m)|=24-10=14>0$) and block weights $(6,2,2)$ up to permutation. The unique block of weight $6$ identifies $\mathcal{B}_{h_k}$ and hence $B_k$; within $\mathcal{B}_{h_k}$, $f_m$ takes value $1$ on all $4$ supports containing $v_m$ and on only $2$ of the $4$ supports containing its partner $v'_m\in B_k$, which uniquely identifies $v_m$; and subtracting $\mathbf{1}_{v_m}$ from $f_m$ recovers $\mathbf{1}_{u_m}$, which uniquely identifies $u_m\in C_1\cup C_2$ since all $17$ core vertex profiles are distinct. Thus all $24$ zero-profiles $Z(e_1),\ldots,Z(e_{24})$ are nonempty and pairwise distinct, and Lemma~\ref{lem:plane-completion}(3) guarantees that the $48$ canonical states $s_{\alpha,\varepsilon}$ are unital and separating on all $65$ vertices of $H_{65}$.
\end{proof}


\subsection{Nonconformality and the missing hub context}\label{sec:nonconformality}

A hypergraph is \emph{conformal} if every clique of its context graph is contained in a hyperedge. The maximal cliques of $\Gamma(H_{65})$ are its $29$ context $4$-cliques and the $13$ uncovered triangles listed in the proof of Theorem~\ref{thm:h65}. Thus $H_{65}$ is nonconformal. Each listed context is nevertheless a complete orthogonal basis, every available $4$-ray basis is listed, and every orthogonal pair belongs to a listed context, as required by Eq.~\eqref{eq:edge-faithful}. The uncovered triangles are orthogonal triples whose fourth basis ray is absent from the selected configuration.

For the hub triangle $\{9,10,11\}$, the missing ray is represented by $q=(1,1,1,-1)^T$. Its omission matters for the state space. The core-state classification gives
\begin{equation}\label{eq:hub-state-sum}
s(9)+s(10)+s(11)=1\qquad(s\in\mathcal{S}(H_{65})).
\end{equation}
Adjoin a new vertex $q_0$ represented by $q$ and the context $\{9,10,11,q_0\}$. Every extended state $\widetilde s$ restricts to an original state, so Eq.~\eqref{eq:hub-state-sum} forces
\begin{equation}\label{eq:hub-extension-zero}
\widetilde s(q_0)=1-\sum_{i=9}^{11}\widetilde s(i)=0.
\end{equation}
Consequently, even completing this single orthogonal triple destroys unitality at the new vertex. Further added contexts cannot remove the implication in Eq.~\eqref{eq:hub-extension-zero} while the original contexts are retained. Completing all orthogonal subsets within the selected ray set would therefore be an additional requirement that this example does not satisfy.

\subsection{Scope of the partition representation}\label{sec:partition-scope}

For either the canonical family or the full state space, the sets $A_i$ of Eq.~\eqref{eq:set-representation} separate all named vertices and turn the listed contexts into partitions, while remaining weaker than an order-determining or compound-event-faithful representation (as already seen for the warm-up $H_{\YO}^{(4)}$ in Appendix~\ref{app:states-hyo4}, where the four non-orthogonal triad atoms $10,11,12,13$ have pairwise disjoint state sets $A_{10},A_{11},A_{12},A_{13}$ and satisfy strict inclusions such as $A_{11}\subsetneq A_{25}$). For $H_{65}$, the classification of its core supports similarly implies
\begin{equation}\label{eq:compound-collapse}
A_9=A_{16}\mathbin{\dot\cup}A_{17},\quad
A_{10}=A_{14}\mathbin{\dot\cup}A_{15},\quad
A_{11}=A_{12}\mathbin{\dot\cup}A_{13}.
\end{equation}
Here $\dot\cup$ denotes disjoint union. Since both $A_{16}$ and $A_{17}$ are nonempty, the first identity gives $A_{16}\subsetneq A_9$. In particular,
\begin{equation}\label{eq:atomic-order-failure}
s(16)\le s(9)\qquad\text{for every }s\in\mathcal{S}(H_{65}).
\end{equation}
Yet $[v_{16}]$ and $[v_9]$ are distinct one-dimensional subspaces, so neither contains the other. Equation~\eqref{eq:atomic-order-failure} holds on the full state space; selecting a different subfamily cannot remove this false Hilbert-space implication. Thus $H_{65}$ fails order determination already on atoms.

Moreover, $v_9=v_{16}-v_{17}$, and the independent orthogonal vectors $v_{16},v_{17}$ span a plane strictly containing $[v_9]$. Equation~\eqref{eq:compound-collapse} nevertheless gives that ray and plane the same set image, because the plane is the compound event obtained by joining vertices $16,17$ in $C_5$. Hence no unital, separating subfamily can make the state-induced representation injective on these compound Hilbert events. The corresponding failure for the warm-up is explicit in Appendix~\ref{app:states-hyo4}.

This distinction also appears at the probability level. Let $q=(1,1,1,-1)^T$, the missing fourth odd-parity Hadamard ray, and define $\Pi_i=v_iv_i^T/\|v_i\|^2$ and $\Pi_q=qq^T/4$. Direct summation gives
\begin{equation}\label{eq:core-projector-sum}
\sum_{i=1}^{17}\Pi_i=4I_4+\Pi_q.
\end{equation}
Every two-valued state has exactly four active core vertices, so every convex mixture of such states gives total core probability $4$. By Eq.~\eqref{eq:core-projector-sum}, quantum preparation along $q$, represented by the unit vector $q/2$, gives total core probability $5$. The existence of separating two-valued states therefore does not imply a classical model for all quantum preparations. Appendix~\ref{app:state-independence} gives the corresponding probability inequalities and an explicit noncontextual quantum state, clarifying the limits of state independence for the fixed four-dimensional realizations.

\paragraph*{Size and completion vertices.}
The $48$ fresh vertices have context degree one; they supply full bases while the $17$-ray core supplies the chromatic obstruction. Deleting one leaves a basis incomplete, and deleting its context leaves a core orthogonal pair uncovered. No minimality is established. Smaller linear real four-dimensional examples with the same properties, fewer auxiliary rays, or no degree-one vertices remain open possibilities.

\appendix

\section{Explicit enumeration of unnormalized vectors, contexts, and non-collinearity certificates}

For direct human and computational verification, this appendix enumerates all unnormalized integer vectors in $\mathbb{Z}^4$ and all four-element contexts for both hypergraphs studied in this paper:
\begin{enumerate}
\item The \emph{canonical linear hypergraph $H_{65}=(V_{65},E_{29})$} (Theorem~\ref{thm:h65}, Appendix~\ref{app:h65}), which has exactly \textbf{$65$ pairwise non-collinear integer vectors} ($17$ core vectors $v_1,\ldots,v_{17}$ and $48$ degree-$1$ completion vectors $v_{18},\ldots,v_{65}$), \textbf{$29$ four-element contexts} ($5$ core contexts $C_1,\ldots,C_5$ and $24$ bridge-completion contexts $E_1,\ldots,E_{24}$), and \textbf{$174$ orthogonality edges} ($29\times\binom{4}{2}=174$).
\item The \emph{non-linear warm-up apex suspension $H_{\YO}^{(4)}=(V_{26},E_{16}^{(4)})$} (Proposition~\ref{prop:warmup}, Appendix~\ref{app:hyo4}), which has exactly \textbf{$26$ pairwise non-collinear integer vectors} ($13$ Yu--Oh rays $v_1,\ldots,v_{13}$, $12$ completion rays $v_{14},\ldots,v_{25}$ in $\R^3\times\{0\}$, and $1$ apex ray $v_{26}$), \textbf{$16$ four-element contexts}, and \textbf{$73$ orthogonality edges} ($16\times\binom{3}{2}=48$ edges in $\R^3\times\{0\}$ plus $25$ apex edges).
\end{enumerate}

\paragraph*{Canonical primitive normal form and human-verifiable non-collinearity.}
Every vector $v_i=(x_{i,1},x_{i,2},x_{i,3},x_{i,4})^T\in\mathbb{Z}^4$ listed below in both $H_{65}$ and $H_{\YO}^{(4)}$ is written in \emph{canonical primitive projective form}:
\begin{equation}\label{eq:canonical-primitive}
\begin{gathered}
\gcd\bigl(|x_{i,1}|,|x_{i,2}|,|x_{i,3}|,|x_{i,4}|\bigr)=1,\\
\text{first nonzero entry of }v_i > 0.
\end{gathered}
\end{equation}
Because every projective ray $[v]\in\PP(\Q^4)$ possesses a \emph{unique} representative in $\mathbb{Z}^4$ satisfying Eq.~\eqref{eq:canonical-primitive}, two vectors $v_i,v_j$ in canonical primitive form are collinear (proportional over $\R$ or $\C$) \emph{if and only if they are identical $4$-tuples} ($v_i=v_j$) and, in particular, have the same squared Euclidean norm $\|v_i\|^2=\|v_j\|^2$.

\subsection{The canonical linear \texorpdfstring{$65$}{65}-vertex, \texorpdfstring{$29$}{29}-context hypergraph \texorpdfstring{$H_{65}$}{H\_65}}\label{app:h65}

Table~\ref{tab:app-h65-vectors} enumerates all $65$ unnormalized integer vectors $v_1,\ldots,v_{65}\in\mathbb{Z}^4$ of $H_{65}$, together with their squared norms $\|v_i\|^2$, their colors $c(i)\in\{0,1,2,3,4\}$ in the strong $5$-coloring of Eq.~\eqref{eq:h65-coordinates}, and the contexts in $E_{29}$ to which each vertex belongs. Table~\ref{tab:app-h65-contexts} enumerates all $29$ contexts $C_1,\ldots,C_5,E_1,\ldots,E_{24}\in E_{29}$, displaying side by side both the vertex indices $\{i,j,k,\ell\}$ and the four mutually orthogonal integer vectors $\{v_i,v_j,v_k,v_\ell\}\subset\mathbb{Z}^4$.

To verify by hand that all $65$ vectors in Table~\ref{tab:app-h65-vectors} are pairwise non-collinear across all $\binom{65}{2}=2{,}080$ pairs, observe that all $65$ tuples satisfy Eq.~\eqref{eq:canonical-primitive} and partition by squared norm $\|v_i\|^2$ into seven small, disjoint norm shells ($4+6+7+12+12+12+12=65$):
\begin{itemize}
\item $\|v_i\|^2=1$ ($4$ vectors): indices $\{1,2,3,4\}$, the $4$ standard coordinate basis vectors;
\item $\|v_i\|^2=2$ ($6$ vectors): indices $\{12,\ldots,17\}$, one vector for each of the $\binom{4}{2}=6$ coordinate-pair supports;
\item $\|v_i\|^2=4$ ($7$ vectors): indices $\{5,\ldots,11\}$, seven distinct $\pm1$ vectors with $x_{i,1}=+1$;
\item $\|v_i\|^2=11$ ($12$ vectors): indices $\{18,20,22,25,26,28,30,33,34,36,39,41\}$, each having a single zero coordinate (three each with $x_1=0$, $x_2=0$, $x_3=0$, $x_4=0$) and distinct sign/entry patterns of $\{0,\pm1,\pm1,\pm3\}$;
\item $\|v_i\|^2=18$ ($12$ vectors): indices $\{43,45,46,49,51,53,54,56,59,61,62,65\}$, twelve distinct permutations and sign patterns of $\{\pm1,\pm2,\pm2,\pm3\}$;
\item $\|v_i\|^2=22$ ($12$ vectors): indices $\{19,21,23,24,27,29,31,32,35,37,38,40\}$, each having a single zero coordinate (three each with $x_1=0$, $x_2=0$, $x_3=0$, $x_4=0$) and distinct patterns of $\{0,\pm2,\pm3,\pm3\}$;
\item $\|v_i\|^2=36$ ($12$ vectors): indices $\{42,44,47,48,50,52,55,57,58,60,63,64\}$, twelve distinct permutations and sign patterns of $\{\pm1,\pm1,\pm3,\pm5\}$.
\end{itemize}
Because no two $4$-tuples within any of the seven norm shells coincide, no two of the $65$ vectors are collinear; quantitatively, the maximum squared cosine over all $\binom{65}{2}=2{,}080$ pairs is $\max_{1\le i<j\le 65}(v_i\cdot v_j)^2/(\|v_i\|^2\|v_j\|^2)=19^2/(11\times 36)=361/396\approx 0.9116<1$, attained by $(v_{41},v_{47})$, whose squared norms are $11$ and $36$ and whose inner product is $19$.

\begin{table*}[t]
\caption{Complete enumeration of the $65$ unnormalized integer vectors $v_i\in\mathbb{Z}^4$ of $H_{65}$, their squared norms $\|v_i\|^2$, their colors $c(i)\in\{0,1,2,3,4\}$, and the contexts in $E_{29}$ containing vertex $i$.}\label{tab:app-h65-vectors}
\scriptsize
\begin{ruledtabular}
\begin{tabular}{rllll|rllll}
$i$ & $v_i\in\mathbb{Z}^4$ & $\|v_i\|^2$ & $c(i)$ & Contexts containing $i$ &
$i$ & $v_i\in\mathbb{Z}^4$ & $\|v_i\|^2$ & $c(i)$ & Contexts containing $i$ \\ \hline
$1$  & $(1,0,0,0)$     & $1$  & $0$ & $C_1,E_3,E_7,E_{11}$ &
$34$ & $(1,-1,0,3)$    & $11$ & $1$ & $E_9=\{3,16,34,35\}$ \\
$2$  & $(0,1,0,0)$     & $1$  & $1$ & $C_1,E_1,E_5,E_{12}$ &
$35$ & $(3,-3,0,-2)$   & $22$ & $3$ & $E_9=\{3,16,34,35\}$ \\
$3$  & $(0,0,1,0)$     & $1$  & $2$ & $C_1,E_2,E_8,E_9$ &
$36$ & $(1,-1,3,0)$    & $11$ & $1$ & $E_{10}=\{4,16,36,37\}$ \\
$4$  & $(0,0,0,1)$     & $1$  & $3$ & $C_1,E_4,E_6,E_{10}$ &
$37$ & $(3,-3,-2,0)$   & $22$ & $2$ & $E_{10}=\{4,16,36,37\}$ \\
$5$  & $(1,1,1,1)$     & $4$  & $0$ & $C_2,E_{13},E_{19},E_{23}$ &
$38$ & $(0,2,3,3)$     & $22$ & $1$ & $E_{11}=\{1,17,38,39\}$ \\
$6$  & $(1,1,-1,-1)$   & $4$  & $1$ & $C_2,E_{15},E_{17},E_{24}$ &
$39$ & $(0,3,-1,-1)$   & $11$ & $3$ & $E_{11}=\{1,17,38,39\}$ \\
$7$  & $(1,-1,1,-1)$   & $4$  & $2$ & $C_2,E_{16},E_{20},E_{21}$ &
$40$ & $(2,0,-3,-3)$   & $22$ & $0$ & $E_{12}=\{2,17,40,41\}$ \\
$8$  & $(1,-1,-1,1)$   & $4$  & $4$ & $C_2,E_{14},E_{18},E_{22}$ &
$41$ & $(3,0,1,1)$     & $11$ & $3$ & $E_{12}=\{2,17,40,41\}$ \\
$9$  & $(1,1,-1,1)$    & $4$  & $0$ & $C_3,C_4$ &
$42$ & $(1,-5,3,1)$    & $36$ & $1$ & $E_{13}=\{5,12,42,43\}$ \\
$10$ & $(1,-1,1,1)$    & $4$  & $1$ & $C_3,C_5$ &
$43$ & $(2,-1,-3,2)$   & $18$ & $2$ & $E_{13}=\{5,12,42,43\}$ \\
$11$ & $(1,-1,-1,-1)$  & $4$  & $3$ & $C_4,C_5$ &
$44$ & $(1,-3,5,1)$    & $36$ & $0$ & $E_{14}=\{8,12,44,45\}$ \\
$12$ & $(1,0,0,-1)$    & $2$  & $3$ & $C_3,E_1,E_2,E_{13},E_{14}$ &
$45$ & $(2,3,1,2)$     & $18$ & $1$ & $E_{14}=\{8,12,44,45\}$ \\
$13$ & $(0,1,1,0)$     & $2$  & $4$ & $C_3,E_3,E_4,E_{15},E_{16}$ &
$46$ & $(1,-2,2,-3)$   & $18$ & $0$ & $E_{15}=\{6,13,46,47\}$ \\
$14$ & $(1,0,1,0)$     & $2$  & $2$ & $C_4,E_5,E_6,E_{17},E_{18}$ &
$47$ & $(5,-1,1,3)$    & $36$ & $2$ & $E_{15}=\{6,13,46,47\}$ \\
$15$ & $(0,1,0,-1)$    & $2$  & $1$ & $C_4,E_7,E_8,E_{19},E_{20}$ &
$48$ & $(3,-1,1,5)$    & $36$ & $0$ & $E_{16}=\{7,13,48,49\}$ \\
$16$ & $(1,1,0,0)$     & $2$  & $0$ & $C_5,E_9,E_{10},E_{21},E_{22}$ &
$49$ & $(3,2,-2,-1)$   & $18$ & $1$ & $E_{16}=\{7,13,48,49\}$ \\
$17$ & $(0,0,1,-1)$    & $2$  & $2$ & $C_5,E_{11},E_{12},E_{23},E_{24}$ &
$50$ & $(1,3,-1,5)$    & $36$ & $0$ & $E_{17}=\{6,14,50,51\}$ \\
$18$ & $(1,0,-3,1)$    & $11$ & $0$ & $E_1=\{2,12,18,19\}$ &
$51$ & $(2,-3,-2,1)$   & $18$ & $3$ & $E_{17}=\{6,14,50,51\}$ \\
$19$ & $(3,0,2,3)$     & $22$ & $2$ & $E_1=\{2,12,18,19\}$ &
$52$ & $(1,5,-1,3)$    & $36$ & $0$ & $E_{18}=\{8,14,52,53\}$ \\
$20$ & $(1,3,0,1)$     & $11$ & $0$ & $E_2=\{3,12,20,21\}$ &
$53$ & $(2,1,-2,-3)$   & $18$ & $1$ & $E_{18}=\{8,14,52,53\}$ \\
$21$ & $(3,-2,0,3)$    & $22$ & $1$ & $E_2=\{3,12,20,21\}$ &
$54$ & $(3,-2,1,-2)$   & $18$ & $2$ & $E_{19}=\{5,15,54,55\}$ \\
$22$ & $(0,1,-1,3)$    & $11$ & $1$ & $E_3=\{1,13,22,23\}$ &
$55$ & $(3,1,-5,1)$    & $36$ & $3$ & $E_{19}=\{5,15,54,55\}$ \\
$23$ & $(0,3,-3,-2)$   & $22$ & $2$ & $E_3=\{1,13,22,23\}$ &
$56$ & $(1,2,3,2)$     & $18$ & $0$ & $E_{20}=\{7,15,56,57\}$ \\
$24$ & $(2,-3,3,0)$    & $22$ & $0$ & $E_4=\{4,13,24,25\}$ &
$57$ & $(5,1,-3,1)$    & $36$ & $3$ & $E_{20}=\{7,15,56,57\}$ \\
$25$ & $(3,1,-1,0)$    & $11$ & $1$ & $E_4=\{4,13,24,25\}$ &
$58$ & $(1,-1,-5,-3)$  & $36$ & $1$ & $E_{21}=\{7,16,58,59\}$ \\
$26$ & $(1,0,-1,-3)$   & $11$ & $0$ & $E_5=\{2,14,26,27\}$ &
$59$ & $(2,-2,-1,3)$   & $18$ & $3$ & $E_{21}=\{7,16,58,59\}$ \\
$27$ & $(3,0,-3,2)$    & $22$ & $3$ & $E_5=\{2,14,26,27\}$ &
$60$ & $(1,-1,-3,-5)$  & $36$ & $1$ & $E_{22}=\{8,16,60,61\}$ \\
$28$ & $(1,-3,-1,0)$   & $11$ & $0$ & $E_6=\{4,14,28,29\}$ &
$61$ & $(2,-2,3,-1)$   & $18$ & $2$ & $E_{22}=\{8,16,60,61\}$ \\
$29$ & $(3,2,-3,0)$    & $22$ & $1$ & $E_6=\{4,14,28,29\}$ &
$62$ & $(1,3,-2,-2)$   & $18$ & $1$ & $E_{23}=\{5,17,62,63\}$ \\
$30$ & $(0,1,3,1)$     & $11$ & $2$ & $E_7=\{1,15,30,31\}$ &
$63$ & $(5,-3,-1,-1)$  & $36$ & $3$ & $E_{23}=\{5,17,62,63\}$ \\
$31$ & $(0,3,-2,3)$    & $22$ & $3$ & $E_7=\{1,15,30,31\}$ &
$64$ & $(3,-5,-1,-1)$  & $36$ & $0$ & $E_{24}=\{6,17,64,65\}$ \\
$32$ & $(2,3,0,3)$     & $22$ & $0$ & $E_8=\{3,15,32,33\}$ &
$65$ & $(3,1,2,2)$     & $18$ & $3$ & $E_{24}=\{6,17,64,65\}$ \\
$33$ & $(3,-1,0,-1)$   & $11$ & $3$ & $E_8=\{3,15,32,33\}$ &
     &                 &      &     &
\end{tabular}
\end{ruledtabular}
\end{table*}

\begin{table*}[t]
\caption{Complete enumeration of the $29$ four-element contexts of $H_{65}=(V_{65},E_{29})$, listing for each context both its vertex indices $\{i,j,k,\ell\}$ and the four mutually orthogonal unnormalized integer vectors $\{v_i,v_j,v_k,v_\ell\}\subset\mathbb{Z}^4$.}\label{tab:app-h65-contexts}
\scriptsize
\begin{ruledtabular}
\begin{tabular}{llllll}
Context & Vertex set & $v_i$ & $v_j$ & $v_k$ & $v_\ell$ \\ \hline
\multicolumn{6}{l}{\emph{Five core $4$-contexts $C_1,\ldots,C_5$ on $V_{\mathrm{core}}=\{1,\ldots,17\}$:}} \\
$C_1$ & $\{1,2,3,4\}$     & $v_1=(1,0,0,0)$   & $v_2=(0,1,0,0)$    & $v_3=(0,0,1,0)$     & $v_4=(0,0,0,1)$ \\
$C_2$ & $\{5,6,7,8\}$     & $v_5=(1,1,1,1)$   & $v_6=(1,1,-1,-1)$  & $v_7=(1,-1,1,-1)$   & $v_8=(1,-1,-1,1)$ \\
$C_3$ & $\{9,10,12,13\}$  & $v_9=(1,1,-1,1)$  & $v_{10}=(1,-1,1,1)$& $v_{12}=(1,0,0,-1)$ & $v_{13}=(0,1,1,0)$ \\
$C_4$ & $\{9,11,14,15\}$  & $v_9=(1,1,-1,1)$  & $v_{11}=(1,-1,-1,-1)$ & $v_{14}=(1,0,1,0)$ & $v_{15}=(0,1,0,-1)$ \\
$C_5$ & $\{10,11,16,17\}$ & $v_{10}=(1,-1,1,1)$ & $v_{11}=(1,-1,-1,-1)$ & $v_{16}=(1,1,0,0)$ & $v_{17}=(0,0,1,-1)$ \\ \hline
\multicolumn{6}{l}{\emph{Twelve coordinate-to-bridge completion contexts $E_1,\ldots,E_{12}$ (between $C_1$ and $B_1\cup B_2\cup B_3$):}} \\
$E_1$  & $\{2,12,18,19\}$ & $v_2=(0,1,0,0)$ & $v_{12}=(1,0,0,-1)$ & $v_{18}=(1,0,-3,1)$  & $v_{19}=(3,0,2,3)$ \\
$E_2$  & $\{3,12,20,21\}$ & $v_3=(0,0,1,0)$ & $v_{12}=(1,0,0,-1)$ & $v_{20}=(1,3,0,1)$   & $v_{21}=(3,-2,0,3)$ \\
$E_3$  & $\{1,13,22,23\}$ & $v_1=(1,0,0,0)$ & $v_{13}=(0,1,1,0)$  & $v_{22}=(0,1,-1,3)$  & $v_{23}=(0,3,-3,-2)$ \\
$E_4$  & $\{4,13,24,25\}$ & $v_4=(0,0,0,1)$ & $v_{13}=(0,1,1,0)$  & $v_{24}=(2,-3,3,0)$  & $v_{25}=(3,1,-1,0)$ \\
$E_5$  & $\{2,14,26,27\}$ & $v_2=(0,1,0,0)$ & $v_{14}=(1,0,1,0)$  & $v_{26}=(1,0,-1,-3)$ & $v_{27}=(3,0,-3,2)$ \\
$E_6$  & $\{4,14,28,29\}$ & $v_4=(0,0,0,1)$ & $v_{14}=(1,0,1,0)$  & $v_{28}=(1,-3,-1,0)$ & $v_{29}=(3,2,-3,0)$ \\
$E_7$  & $\{1,15,30,31\}$ & $v_1=(1,0,0,0)$ & $v_{15}=(0,1,0,-1)$ & $v_{30}=(0,1,3,1)$   & $v_{31}=(0,3,-2,3)$ \\
$E_8$  & $\{3,15,32,33\}$ & $v_3=(0,0,1,0)$ & $v_{15}=(0,1,0,-1)$ & $v_{32}=(2,3,0,3)$   & $v_{33}=(3,-1,0,-1)$ \\
$E_9$  & $\{3,16,34,35\}$ & $v_3=(0,0,1,0)$ & $v_{16}=(1,1,0,0)$  & $v_{34}=(1,-1,0,3)$  & $v_{35}=(3,-3,0,-2)$ \\
$E_{10}$ & $\{4,16,36,37\}$ & $v_4=(0,0,0,1)$ & $v_{16}=(1,1,0,0)$  & $v_{36}=(1,-1,3,0)$  & $v_{37}=(3,-3,-2,0)$ \\
$E_{11}$ & $\{1,17,38,39\}$ & $v_1=(1,0,0,0)$ & $v_{17}=(0,0,1,-1)$ & $v_{38}=(0,2,3,3)$   & $v_{39}=(0,3,-1,-1)$ \\
$E_{12}$ & $\{2,17,40,41\}$ & $v_2=(0,1,0,0)$ & $v_{17}=(0,0,1,-1)$ & $v_{40}=(2,0,-3,-3)$ & $v_{41}=(3,0,1,1)$ \\ \hline
\multicolumn{6}{l}{\emph{Twelve Hadamard-to-bridge completion contexts $E_{13},\ldots,E_{24}$ (between $C_2$ and $B_1\cup B_2\cup B_3$):}} \\
$E_{13}$ & $\{5,12,42,43\}$ & $v_5=(1,1,1,1)$    & $v_{12}=(1,0,0,-1)$ & $v_{42}=(1,-5,3,1)$   & $v_{43}=(2,-1,-3,2)$ \\
$E_{14}$ & $\{8,12,44,45\}$ & $v_8=(1,-1,-1,1)$  & $v_{12}=(1,0,0,-1)$ & $v_{44}=(1,-3,5,1)$   & $v_{45}=(2,3,1,2)$ \\
$E_{15}$ & $\{6,13,46,47\}$ & $v_6=(1,1,-1,-1)$  & $v_{13}=(0,1,1,0)$  & $v_{46}=(1,-2,2,-3)$  & $v_{47}=(5,-1,1,3)$ \\
$E_{16}$ & $\{7,13,48,49\}$ & $v_7=(1,-1,1,-1)$  & $v_{13}=(0,1,1,0)$  & $v_{48}=(3,-1,1,5)$   & $v_{49}=(3,2,-2,-1)$ \\
$E_{17}$ & $\{6,14,50,51\}$ & $v_6=(1,1,-1,-1)$  & $v_{14}=(1,0,1,0)$  & $v_{50}=(1,3,-1,5)$   & $v_{51}=(2,-3,-2,1)$ \\
$E_{18}$ & $\{8,14,52,53\}$ & $v_8=(1,-1,-1,1)$  & $v_{14}=(1,0,1,0)$  & $v_{52}=(1,5,-1,3)$   & $v_{53}=(2,1,-2,-3)$ \\
$E_{19}$ & $\{5,15,54,55\}$ & $v_5=(1,1,1,1)$    & $v_{15}=(0,1,0,-1)$ & $v_{54}=(3,-2,1,-2)$  & $v_{55}=(3,1,-5,1)$ \\
$E_{20}$ & $\{7,15,56,57\}$ & $v_7=(1,-1,1,-1)$  & $v_{15}=(0,1,0,-1)$ & $v_{56}=(1,2,3,2)$    & $v_{57}=(5,1,-3,1)$ \\
$E_{21}$ & $\{7,16,58,59\}$ & $v_7=(1,-1,1,-1)$  & $v_{16}=(1,1,0,0)$  & $v_{58}=(1,-1,-5,-3)$ & $v_{59}=(2,-2,-1,3)$ \\
$E_{22}$ & $\{8,16,60,61\}$ & $v_8=(1,-1,-1,1)$  & $v_{16}=(1,1,0,0)$  & $v_{60}=(1,-1,-3,-5)$ & $v_{61}=(2,-2,3,-1)$ \\
$E_{23}$ & $\{5,17,62,63\}$ & $v_5=(1,1,1,1)$    & $v_{17}=(0,0,1,-1)$ & $v_{62}=(1,3,-2,-2)$  & $v_{63}=(5,-3,-1,-1)$ \\
$E_{24}$ & $\{6,17,64,65\}$ & $v_6=(1,1,-1,-1)$  & $v_{17}=(0,0,1,-1)$ & $v_{64}=(3,-5,-1,-1)$ & $v_{65}=(3,1,2,2)$
\end{tabular}
\end{ruledtabular}
\end{table*}

\subsection{The non-linear warm-up \texorpdfstring{$26$}{26}-vertex, \texorpdfstring{$16$}{16}-context apex suspension \texorpdfstring{$H_{\YO}^{(4)}$}{H\_YO\^(4)}}\label{app:hyo4}

For completeness, the $26$ unnormalized integer vectors $v_1,\ldots,v_{26}\in\mathbb{Z}^4$ of the non-linear warm-up hypergraph $H_{\YO}^{(4)}=(V_{26},E_{16}^{(4)})$ in Proposition~\ref{prop:warmup} (written in canonical primitive form satisfying Eq.~\eqref{eq:canonical-primitive}) are:
\begin{equation}\label{eq:hyo4-vectors}
\begin{aligned}
v_1&=(0,1,1,0),& v_2&=(1,0,0,0),\\
v_3&=(0,1,-1,0),& v_4&=(1,1,0,0),\\
v_5&=(0,0,1,0),& v_6&=(1,-1,0,0),\\
v_7&=(1,0,1,0),& v_8&=(0,1,0,0),\\
v_9&=(1,0,-1,0),& v_{10}&=(1,-1,1,0),\\
v_{11}&=(1,-1,-1,0),& v_{12}&=(1,1,1,0),\\
v_{13}&=(1,1,-1,0),& v_{14}&=(2,1,1,0),\\
v_{15}&=(1,2,-1,0),& v_{16}&=(1,-2,1,0),\\
v_{17}&=(1,1,-2,0),& v_{18}&=(1,-1,-2,0),\\
v_{19}&=(2,1,-1,0),& v_{20}&=(2,-1,1,0),\\
v_{21}&=(1,-2,-1,0),& v_{22}&=(1,1,2,0),\\
v_{23}&=(1,2,1,0),& v_{24}&=(1,-1,2,0),\\
v_{25}&=(2,-1,-1,0),& v_{26}&=(0,0,0,1).
\end{aligned}
\end{equation}
These $26$ canonical primitive vectors partition by squared norm $\|v_i\|^2$ into four disjoint norm shells ($4+6+4+12=26$):
\begin{itemize}
\item $\|v_i\|^2=1$ ($4$ vectors): $\{v_2,v_5,v_8,v_{26}\}$, the four coordinate basis vectors of $\R^4$;
\item $\|v_i\|^2=2$ ($6$ vectors): $\{v_1,v_3,v_4,v_6,v_7,v_9\}$, the six binary rays in $\R^3\times\{0\}$;
\item $\|v_i\|^2=3$ ($4$ vectors): $\{v_{10},v_{11},v_{12},v_{13}\}$, the four triad rays $(1,\pm1,\pm1,0)^T$;
\item $\|v_i\|^2=6$ ($12$ vectors): $\{v_{14},\ldots,v_{25}\}$, all $3\times 4=12$ canonical primitive permutations and sign patterns of $(\pm2,\pm1,\pm1,0)^T$.
\end{itemize}
Since all $26$ canonical $4$-tuples are distinct, all $\binom{26}{2}=325$ pairs of vectors are strictly non-collinear, with maximum squared cosine $\max_{1\le i<j\le 26}(v_i\cdot v_j)^2/(\|v_i\|^2\|v_j\|^2)=4^2/(3\times 6)=8/9\approx 0.8889<1$. Table~\ref{tab:app-hyo4-contexts} lists all $16$ four-element contexts of $E_{16}^{(4)}$ side by side with their four mutually orthogonal integer vectors.

\begin{table*}[t]
\caption{Complete enumeration of the $16$ four-element contexts of the non-linear warm-up apex suspension $H_{\YO}^{(4)}=(V_{26},E_{16}^{(4)})$, listing for each context both its vertex indices $\{i,j,k,26\}$ and the four mutually orthogonal unnormalized integer vectors $\{v_i,v_j,v_k,v_{26}\}\subset\mathbb{Z}^4$.}\label{tab:app-hyo4-contexts}
\scriptsize
\begin{ruledtabular}
\begin{tabular}{llllll}
Context & Vertex set & $v_i$ & $v_j$ & $v_k$ & $v_{26}$ \\ \hline
$e_1^{(4)}$  & $\{1,2,3,26\}$   & $v_1=(0,1,1,0)$  & $v_2=(1,0,0,0)$    & $v_3=(0,1,-1,0)$    & $v_{26}=(0,0,0,1)$ \\
$e_2^{(4)}$  & $\{7,8,9,26\}$   & $v_7=(1,0,1,0)$  & $v_8=(0,1,0,0)$    & $v_9=(1,0,-1,0)$    & $v_{26}=(0,0,0,1)$ \\
$e_3^{(4)}$  & $\{4,5,6,26\}$   & $v_4=(1,1,0,0)$  & $v_5=(0,0,1,0)$    & $v_6=(1,-1,0,0)$    & $v_{26}=(0,0,0,1)$ \\
$e_4^{(4)}$  & $\{2,5,8,26\}$   & $v_2=(1,0,0,0)$  & $v_5=(0,0,1,0)$    & $v_8=(0,1,0,0)$     & $v_{26}=(0,0,0,1)$ \\ \hline
$e_5^{(4)}$  & $\{3,12,25,26\}$ & $v_3=(0,1,-1,0)$ & $v_{12}=(1,1,1,0)$ & $v_{25}=(2,-1,-1,0)$& $v_{26}=(0,0,0,1)$ \\
$e_6^{(4)}$  & $\{9,12,16,26\}$ & $v_9=(1,0,-1,0)$ & $v_{12}=(1,1,1,0)$ & $v_{16}=(1,-2,1,0)$ & $v_{26}=(0,0,0,1)$ \\
$e_7^{(4)}$  & $\{6,12,17,26\}$ & $v_6=(1,-1,0,0)$ & $v_{12}=(1,1,1,0)$ & $v_{17}=(1,1,-2,0)$ & $v_{26}=(0,0,0,1)$ \\
$e_8^{(4)}$  & $\{3,11,14,26\}$ & $v_3=(0,1,-1,0)$ & $v_{11}=(1,-1,-1,0)$&$v_{14}=(2,1,1,0)$  & $v_{26}=(0,0,0,1)$ \\
$e_9^{(4)}$  & $\{7,11,15,26\}$ & $v_7=(1,0,1,0)$  & $v_{11}=(1,-1,-1,0)$&$v_{15}=(1,2,-1,0)$ & $v_{26}=(0,0,0,1)$ \\
$e_{10}^{(4)}$ & $\{4,11,24,26\}$ & $v_4=(1,1,0,0)$  & $v_{11}=(1,-1,-1,0)$&$v_{24}=(1,-1,2,0)$ & $v_{26}=(0,0,0,1)$ \\
$e_{11}^{(4)}$ & $\{9,10,23,26\}$ & $v_9=(1,0,-1,0)$ & $v_{10}=(1,-1,1,0)$ & $v_{23}=(1,2,1,0)$  & $v_{26}=(0,0,0,1)$ \\
$e_{12}^{(4)}$ & $\{1,10,19,26\}$ & $v_1=(0,1,1,0)$  & $v_{10}=(1,-1,1,0)$ & $v_{19}=(2,1,-1,0)$ & $v_{26}=(0,0,0,1)$ \\
$e_{13}^{(4)}$ & $\{4,10,18,26\}$ & $v_4=(1,1,0,0)$  & $v_{10}=(1,-1,1,0)$ & $v_{18}=(1,-1,-2,0)$& $v_{26}=(0,0,0,1)$ \\
$e_{14}^{(4)}$ & $\{1,13,20,26\}$ & $v_1=(0,1,1,0)$  & $v_{13}=(1,1,-1,0)$ & $v_{20}=(2,-1,1,0)$ & $v_{26}=(0,0,0,1)$ \\
$e_{15}^{(4)}$ & $\{7,13,21,26\}$ & $v_7=(1,0,1,0)$  & $v_{13}=(1,1,-1,0)$ & $v_{21}=(1,-2,-1,0)$& $v_{26}=(0,0,0,1)$ \\
$e_{16}^{(4)}$ & $\{6,13,22,26\}$ & $v_6=(1,-1,0,0)$ & $v_{13}=(1,1,-1,0)$ & $v_{22}=(1,1,2,0)$  & $v_{26}=(0,0,0,1)$
\end{tabular}
\end{ruledtabular}
\end{table*}

\section{Explicit enumeration of two-valued states}

This appendix provides the complete enumeration of all two-valued states on both hypergraphs studied in the main text: the $25$ two-valued states of the non-linear warm-up hypergraph $H_{\YO}^{(4)}$ (Appendix~\ref{app:states-hyo4}), and the complete branch-by-branch parametrization of all $24\times 2^{14}=393{,}216$ two-valued states---together with the $48$ canonical unital and separating states $s_{\alpha,\varepsilon}$ and the $65$ vertex state-profiles---of the canonical linear hypergraph $H_{65}$ (Appendix~\ref{app:states-h65}).

\subsection{Complete enumeration of the \texorpdfstring{$25$}{25} two-valued states on \texorpdfstring{$H_{\YO}^{(4)}$}{H\_YO\^(4)}}\label{app:states-hyo4}

Table~\ref{tab:app-hyo4-states} lists all $25$ two-valued states $s_1,\ldots,s_{25}:V_{26}\to\{0,1\}$ on $H_{\YO}^{(4)}=(V_{26},E_{16}^{(4)})$. States $s_1,\ldots,s_{24}$ have $s_r(26)=0$ and correspond to the $24$ two-valued states of the completed Yu--Oh hypergraph $H_{\YO}^{+}$ in $\R^3$ (grouped into $12$ pairs of states sharing the same support on the three coordinate contexts $\{1,2,3\}$, $\{4,5,6\}$, $\{7,8,9\}$), while $s_{25}$ is the unique apex state with support $\{26\}$. Table~\ref{tab:app-hyo4-profiles} displays the transposed $26\times 25$ vertex-by-state binary matrix, listing for each vertex $i\in\{1,\ldots,26\}$ its state weight $\sum_{r=1}^{25}s_r(i)\in\{1,3,8,13\}$, the set of states $A_i=\{s_r : s_r(i)=1\}$, and its $25$-bit profile $(s_1(i),\ldots,s_{25}(i))$.

Inspection of Table~\ref{tab:app-hyo4-profiles} highlights three structural properties of the state space $\mathcal{S}(H_{\YO}^{(4)})$:
\begin{enumerate}
\item \emph{Unitality and atom separation:} Every row of Table~\ref{tab:app-hyo4-profiles} has positive weight ($1,3,8,$ or $13$) and all $26$ binary rows are distinct ($A_i\ne A_j$ for all $1\le i<j\le 26$), so $\mathcal{S}(H_{\YO}^{(4)})$ is nonempty, unital, and separating on atoms.
\item \emph{Exclusion of simultaneous $1$s on the non-orthogonal triad atoms $\{10,11,12,13\}$:} Although the four Yu--Oh triad rays $v_{10},v_{11},v_{12},v_{13}$ are pairwise non-orthogonal ($v_i\cdot v_j=\pm1\ne0$ for $10\le i<j\le 13$), their active state sets in Table~\ref{tab:app-hyo4-profiles} are pairwise disjoint:
\begin{equation}\label{eq:hyo4-triad-disjoint}
\begin{aligned}
A_{10}&=\{s_{13},s_{19},s_{23}\},&
A_{11}&=\{s_{5},s_{7},s_{15}\},\\
A_{12}&=\{s_{1},s_{3},s_{9}\},&
A_{13}&=\{s_{11},s_{17},s_{21}\}.
\end{aligned}
\end{equation}
Consequently, $\sum_{i=10}^{13}s(i)\le1$ holds for every $s\in\mathcal{S}(H_{\YO}^{(4)})$, and no two-valued state assigns $(1,1)$ to any of the $\binom{4}{2}=6$ non-orthogonal pairs $\{10,11\}$, $\{10,12\}$, $\{10,13\}$, $\{11,12\}$, $\{11,13\}$, $\{12,13\}$.
\item \emph{Failure of order-determinacy:} Even though $A_i\ne A_j$ for all $i\ne j$, the state space $\mathcal{S}(H_{\YO}^{(4)})$ is not order-determining on atoms or compound propositions~\cite{svozil-tkadlec}. In particular, each triad atom $t\in\{10,11,12,13\}$ (\emph{weight $3$}) is strictly dominated across all $25$ states ($A_t\subsetneq A_w$, i.e., $s(t)\le s(w)$ for all $s\in\mathcal{S}(H_{\YO}^{(4)})$) by three non-orthogonal completion atoms $w\in\{14,\ldots,25\}$ (\emph{weight $13$}):
\begin{equation}\label{eq:hyo4-inclusions}
\begin{aligned}
A_{10}&\subsetneq A_{16},\,A_{20},\,A_{24},&
A_{11}&\subsetneq A_{18},\,A_{21},\,A_{25},\\
A_{12}&\subsetneq A_{14},\,A_{22},\,A_{23},&
A_{13}&\subsetneq A_{15},\,A_{17},\,A_{19},
\end{aligned}
\end{equation}
together with compound-event inclusions such as $A_{12}\subsetneq A_3\mathbin{\dot\cup}A_{14}$, $A_{11}\subsetneq A_3\mathbin{\dot\cup}A_{25}$, and their unions with the apex set $A_{26}=\{s_{25}\}$.
\end{enumerate}

\begin{table*}[t]
\caption{Complete enumeration of the $25$ two-valued states $s_1,\ldots,s_{25}:V_{26}\to\{0,1\}$ on the non-linear warm-up hypergraph $H_{\YO}^{(4)}=(V_{26},E_{16}^{(4)})$, showing for each state its support on the $13$-ray Yu--Oh core $\{1,\ldots,13\}$, its forced support on $\{14,\ldots,26\}$, its total support size $|S_r|$, and its $26$-bit binary vector $(s_r(1\ldots 13)\,s_r(14\ldots 25)\,s_r(26))$.}\label{tab:app-hyo4-states}
\scriptsize
\begin{ruledtabular}
\begin{tabular}{llllc}
State & Core support $S_r\cap\{1,\ldots,13\}$ & Completion/apex support $S_r\cap\{14,\ldots,26\}$ & $|S_r|$ & Binary vector $(s_r(1\ldots 13)\,s_r(14\ldots 25)\,s_r(26))$ \\ \hline
$s_{1}$ & $\{1,4,8,12\}$ & $\{14,15,21,22,23\}$ & $9$ & \texttt{1001000100010\,110000011100\,0} \\
$s_{2}$ & $\{1,4,8\}$ & $\{14,15,16,17,21,22,23,25\}$ & $11$ & \texttt{1001000100000\,111100011101\,0} \\
$s_{3}$ & $\{1,5,7,12\}$ & $\{14,18,22,23,24\}$ & $9$ & \texttt{1000101000010\,100010001110\,0} \\
$s_{4}$ & $\{1,5,7\}$ & $\{14,16,17,18,22,23,24,25\}$ & $11$ & \texttt{1000101000000\,101110001111\,0} \\
$s_{5}$ & $\{1,5,9,11\}$ & $\{17,18,21,22,25\}$ & $9$ & \texttt{1000100010100\,000110011001\,0} \\
$s_{6}$ & $\{1,5,9\}$ & $\{14,15,17,18,21,22,24,25\}$ & $11$ & \texttt{1000100010000\,110110011011\,0} \\
$s_{7}$ & $\{1,6,8,11\}$ & $\{16,18,21,23,25\}$ & $9$ & \texttt{1000010100100\,001010010101\,0} \\
$s_{8}$ & $\{1,6,8\}$ & $\{14,15,16,18,21,23,24,25\}$ & $11$ & \texttt{1000010100000\,111010010111\,0} \\
$s_{9}$ & $\{2,4,7,12\}$ & $\{14,19,20,22,23\}$ & $9$ & \texttt{0101001000010\,100001101100\,0} \\
$s_{10}$ & $\{2,4,7\}$ & $\{14,16,17,19,20,22,23,25\}$ & $11$ & \texttt{0101001000000\,101101101101\,0} \\
$s_{11}$ & $\{2,4,9,13\}$ & $\{14,15,17,19,25\}$ & $9$ & \texttt{0101000010001\,110101000001\,0} \\
$s_{12}$ & $\{2,4,9\}$ & $\{14,15,17,19,20,21,22,25\}$ & $11$ & \texttt{0101000010000\,110101111001\,0} \\
$s_{13}$ & $\{2,6,7,10\}$ & $\{14,16,20,24,25\}$ & $9$ & \texttt{0100011001000\,101000100011\,0} \\
$s_{14}$ & $\{2,6,7\}$ & $\{14,16,18,19,20,23,24,25\}$ & $11$ & \texttt{0100011000000\,101011100111\,0} \\
$s_{15}$ & $\{2,6,9,11\}$ & $\{18,19,20,21,25\}$ & $9$ & \texttt{0100010010100\,000011110001\,0} \\
$s_{16}$ & $\{2,6,9\}$ & $\{14,15,18,19,20,21,24,25\}$ & $11$ & \texttt{0100010010000\,110011110011\,0} \\
$s_{17}$ & $\{3,4,8,13\}$ & $\{15,16,17,19,23\}$ & $9$ & \texttt{0011000100001\,011101000100\,0} \\
$s_{18}$ & $\{3,4,8\}$ & $\{15,16,17,19,20,21,22,23\}$ & $11$ & \texttt{0011000100000\,011101111100\,0} \\
$s_{19}$ & $\{3,5,7,10\}$ & $\{16,17,20,22,24\}$ & $9$ & \texttt{0010101001000\,001100101010\,0} \\
$s_{20}$ & $\{3,5,7\}$ & $\{16,17,18,19,20,22,23,24\}$ & $11$ & \texttt{0010101000000\,001111101110\,0} \\
$s_{21}$ & $\{3,5,9,13\}$ & $\{15,17,18,19,24\}$ & $9$ & \texttt{0010100010001\,010111000010\,0} \\
$s_{22}$ & $\{3,5,9\}$ & $\{15,17,18,19,20,21,22,24\}$ & $11$ & \texttt{0010100010000\,010111111010\,0} \\
$s_{23}$ & $\{3,6,8,10\}$ & $\{15,16,20,21,24\}$ & $9$ & \texttt{0010010101000\,011000110010\,0} \\
$s_{24}$ & $\{3,6,8\}$ & $\{15,16,18,19,20,21,23,24\}$ & $11$ & \texttt{0010010100000\,011011110110\,0} \\ \hline
$s_{25}$ & $\emptyset$ & $\{26\}$ & $1$ & \texttt{0000000000000\,000000000000\,1}
\end{tabular}
\end{ruledtabular}
\end{table*}

\begin{table*}[t]
\caption{Vertex state profiles across the $25$ two-valued states $s_1,\ldots,s_{25}$ of $H_{\YO}^{(4)}$, showing for each vertex $i\in\{1,\ldots,26\}$ its state weight $\sum_{r=1}^{25}s_r(i)$, the active state indices $\{r : s_r(i)=1\}$, and the $25$-bit profile $(s_{1\ldots 12}(i)\,s_{13\ldots 24}(i)\,s_{25}(i))$. All $26$ rows are nonzero and pairwise distinct.}\label{tab:app-hyo4-profiles}
\scriptsize
\begin{ruledtabular}
\begin{tabular}{rrlc}
Vertex $i$ & Weight $\sum_r s_r(i)$ & Active states $\{r\in\{1,\ldots,25\} : s_r(i)=1\}$ & $25$-bit profile $(s_{1\ldots 12}(i)\,s_{13\ldots 24}(i)\,s_{25}(i))$ \\ \hline
$1$ & $8$ & $\{1,2,3,4,5,6,7,8\}$ & \texttt{111111110000\,000000000000\,0} \\
$2$ & $8$ & $\{9,10,11,12,13,14,15,16\}$ & \texttt{000000001111\,111100000000\,0} \\
$3$ & $8$ & $\{17,18,19,20,21,22,23,24\}$ & \texttt{000000000000\,000011111111\,0} \\
$4$ & $8$ & $\{1,2,9,10,11,12,17,18\}$ & \texttt{110000001111\,000011000000\,0} \\
$5$ & $8$ & $\{3,4,5,6,19,20,21,22\}$ & \texttt{001111000000\,000000111100\,0} \\
$6$ & $8$ & $\{7,8,13,14,15,16,23,24\}$ & \texttt{000000110000\,111100000011\,0} \\
$7$ & $8$ & $\{3,4,9,10,13,14,19,20\}$ & \texttt{001100001100\,110000110000\,0} \\
$8$ & $8$ & $\{1,2,7,8,17,18,23,24\}$ & \texttt{110000110000\,000011000011\,0} \\
$9$ & $8$ & $\{5,6,11,12,15,16,21,22\}$ & \texttt{000011000011\,001100001100\,0} \\ \hline
$10$ & $3$ & $\{13,19,23\}$ & \texttt{000000000000\,100000100010\,0} \\
$11$ & $3$ & $\{5,7,15\}$ & \texttt{000010100000\,001000000000\,0} \\
$12$ & $3$ & $\{1,3,9\}$ & \texttt{101000001000\,000000000000\,0} \\
$13$ & $3$ & $\{11,17,21\}$ & \texttt{000000000010\,000010001000\,0} \\ \hline
$14$ & $13$ & $\{1,2,3,4,6,8,9,10,11,12,13,14,16\}$ & \texttt{111101011111\,110100000000\,0} \\
$15$ & $13$ & $\{1,2,6,8,11,12,16,17,18,21,22,23,24\}$ & \texttt{110001010011\,000111001111\,0} \\
$16$ & $13$ & $\{2,4,7,8,10,13,14,17,18,19,20,23,24\}$ & \texttt{010100110100\,110011110011\,0} \\
$17$ & $13$ & $\{2,4,5,6,10,11,12,17,18,19,20,21,22\}$ & \texttt{010111000111\,000011111100\,0} \\
$18$ & $13$ & $\{3,4,5,6,7,8,14,15,16,20,21,22,24\}$ & \texttt{001111110000\,011100011101\,0} \\
$19$ & $13$ & $\{9,10,11,12,14,15,16,17,18,20,21,22,24\}$ & \texttt{000000001111\,011111011101\,0} \\
$20$ & $13$ & $\{9,10,12,13,14,15,16,18,19,20,22,23,24\}$ & \texttt{000000001101\,111101110111\,0} \\
$21$ & $13$ & $\{1,2,5,6,7,8,12,15,16,18,22,23,24\}$ & \texttt{110011110001\,001101000111\,0} \\
$22$ & $13$ & $\{1,2,3,4,5,6,9,10,12,18,19,20,22\}$ & \texttt{111111001101\,000001110100\,0} \\
$23$ & $13$ & $\{1,2,3,4,7,8,9,10,14,17,18,20,24\}$ & \texttt{111100111100\,010011010001\,0} \\
$24$ & $13$ & $\{3,4,6,8,13,14,16,19,20,21,22,23,24\}$ & \texttt{001101010000\,110100111111\,0} \\
$25$ & $13$ & $\{2,4,5,6,7,8,10,11,12,13,14,15,16\}$ & \texttt{010111110111\,111100000000\,0} \\ \hline
$26$ & $1$ & $\{25\}$ & \texttt{000000000000\,000000000000\,1}
\end{tabular}
\end{ruledtabular}
\end{table*}

\subsection{Complete parametrization of all \texorpdfstring{$393{,}216$}{393,216} two-valued states and the \texorpdfstring{$48$}{48} canonical separating states on \texorpdfstring{$H_{65}$}{H\_65}}\label{app:states-h65}

Every two-valued state $s\in\mathcal{S}(H_{65})$ on the canonical linear hypergraph $H_{65}=(V_{65},E_{29})$ is uniquely indexed by a pair $(\alpha,\bm{\varepsilon})$ consisting of:
\begin{enumerate}
\item a branch index $\alpha\in\{1,\ldots,24\}$ selecting one of the $24$ four-element core state supports $S_\alpha\subset V_{\mathrm{core}}=\{1,\ldots,17\}$; and
\item a $14$-bit binary vector $\bm{\varepsilon}=(\varepsilon_m)_{m\in \mathcal{F}_\alpha}\in\{0,1\}^{14}$ indexed by the $14$ \emph{free completion contexts}
\begin{equation*}
\mathcal{F}_\alpha=\bigl\{m\in\{1,\ldots,24\} : S_\alpha\cap\{u_m,v_m\}=\emptyset\bigr\},
\end{equation*}
whose complement is the set of $10$ \emph{hit completion contexts} $\mathcal{H}_\alpha=\{1,\ldots,24\}\setminus \mathcal{F}_\alpha$.
\end{enumerate}
Explicitly, the state $s_{\alpha,\bm{\varepsilon}}:V_{65}\to\{0,1\}$ takes values
\begin{equation}\label{eq:full-state-param}
\begin{aligned}
s_{\alpha,\bm{\varepsilon}}(i)&=\mathbf{1}_{i\in S_\alpha}\qquad(1\le i\le 17),\\
\bigl(s_{\alpha,\bm{\varepsilon}}&(2m+16),\,s_{\alpha,\bm{\varepsilon}}(2m+17)\bigr)\\
&=\begin{cases}
(0,0), & m\in \mathcal{H}_\alpha,\\
(1-\varepsilon_m,\,\varepsilon_m), & m\in \mathcal{F}_\alpha.
\end{cases}
\end{aligned}
\end{equation}
Because $|\mathcal{F}_\alpha|=14$ for every $\alpha\in\{1,\ldots,24\}$, each branch $\alpha$ yields $2^{14}=16{,}384$ states $s_{\alpha,\bm{\varepsilon}}$, giving $|\mathcal{S}(H_{65})|=24\times 2^{14}=393{,}216$ two-valued states in total. Setting all $14$ free bits equal to $\varepsilon\in\{0,1\}$ ($\varepsilon_m=\varepsilon$ for all $m\in \mathcal{F}_\alpha$) yields the $48$ \emph{canonical states} $s_{\alpha,0},s_{\alpha,1}$ of Eq.~\eqref{eq:canonical-states}, which have support size $|S_\alpha|+14=18$:
\begin{equation}\label{eq:canonical-state-supports}
\begin{aligned}
\mathrm{supp}(s_{\alpha,0}) &= S_\alpha \cup \bigl\{2m+16 : m\in \mathcal{F}_\alpha\bigr\},\\
\mathrm{supp}(s_{\alpha,1}) &= S_\alpha \cup \bigl\{2m+17 : m\in \mathcal{F}_\alpha\bigr\}.
\end{aligned}
\end{equation}

Table~\ref{tab:app-h65-state-branches} enumerates all $24$ branches $\alpha=1,\ldots,24$, specifying $S_\alpha$, the $10$ hit context indices $\mathcal{H}_\alpha\subset\{1,\ldots,24\}$, the $14$ free context indices $\mathcal{F}_\alpha\subset\{1,\ldots,24\}$, and the $24$-bit free-context indicator string. Table~\ref{tab:app-h65-vertex-profiles} lists the resulting state profiles and weights for all $65$ vertices of $H_{65}$ across both the $48$ canonical states $\{s_{\alpha,\varepsilon}\}$ and the full space of $393{,}216$ states: across the $48$ canonical states $(s_{1,0},s_{1,1},\ldots,s_{24,0},s_{24,1})$, each `1` in row $i$ of Table~\ref{tab:app-h65-vertex-profiles} expands to `(1,1)` for core vertices $i\in\{1,\ldots,17\}$, to `(1,0)` for first completion vertices $w_{m,1}=2m+16$, and to `(0,1)` for second completion vertices $w_{m,2}=2m+17$ (and each `0` expands to `(0,0)`), making all $65$ canonical $48$-bit vertex profiles nonzero and pairwise distinct by inspection.

\begin{table*}[t]
\caption{Complete branch-by-branch parametrization of the $24\times 2^{14}=393{,}216$ two-valued states $s_{\alpha,\bm{\varepsilon}}$ and the $48$ canonical states $s_{\alpha,0},s_{\alpha,1}$ on $H_{65}$ via Eqs.~\eqref{eq:full-state-param}--\eqref{eq:canonical-state-supports}. For each $\alpha\in\{1,\ldots,24\}$, the table lists the $4$-element core support $S_\alpha$, the $10$ hit completion context indices $m\in \mathcal{H}_\alpha$ where $(s(2m+16),s(2m+17))=(0,0)$ is forced on $E_m$, the $14$ free completion context indices $m\in \mathcal{F}_\alpha$ where $(s(2m+16),s(2m+17))=(1-\varepsilon_m,\varepsilon_m)$ with $\varepsilon_m\in\{0,1\}$ on $E_m$, and the $24$-bit indicator string $(\mathbf{1}_{1\in \mathcal{F}_\alpha}\ldots\mathbf{1}_{12\in \mathcal{F}_\alpha}\,\mathbf{1}_{13\in \mathcal{F}_\alpha}\ldots\mathbf{1}_{24\in \mathcal{F}_\alpha})$.}\label{tab:app-h65-state-branches}
\scriptsize
\begin{ruledtabular}
\begin{tabular}{llllc}
Branch & Core support $S_\alpha$ & Hit contexts $\mathcal{H}_\alpha$ (forced $(0,0)$) & Free contexts $\mathcal{F}_\alpha$ ($2^{14}$ choices) & Free indicator $m=1\ldots 24$ \\ \hline
$S_{1}$ & $\{1,5,9,16\}$ & $\{3,7,9,10,11,13,19,21,22,23\}$ & $\{1,2,4,5,6,8,12,14,15,16,17,18,20,24\}$ & \texttt{110111010001\,011111010001} \\
$S_{2}$ & $\{1,5,10,14\}$ & $\{3,5,6,7,11,13,17,18,19,23\}$ & $\{1,2,4,8,9,10,12,14,15,16,20,21,22,24\}$ & \texttt{110100011101\,011100011101} \\
$S_{3}$ & $\{1,6,9,16\}$ & $\{3,7,9,10,11,15,17,21,22,24\}$ & $\{1,2,4,5,6,8,12,13,14,16,18,19,20,23\}$ & \texttt{110111010001\,110101110010} \\
$S_{4}$ & $\{1,6,11,12\}$ & $\{1,2,3,7,11,13,14,15,17,24\}$ & $\{4,5,6,8,9,10,12,16,18,19,20,21,22,23\}$ & \texttt{000111011101\,000101111110} \\
$S_{5}$ & $\{1,7,10,14\}$ & $\{3,5,6,7,11,16,17,18,20,21\}$ & $\{1,2,4,8,9,10,12,13,14,15,19,22,23,24\}$ & \texttt{110100011101\,111000100111} \\
$S_{6}$ & $\{1,7,11,12\}$ & $\{1,2,3,7,11,13,14,16,20,21\}$ & $\{4,5,6,8,9,10,12,15,17,18,19,22,23,24\}$ & \texttt{000111011101\,001011100111} \\
$S_{7}$ & $\{2,5,9,16\}$ & $\{1,5,9,10,12,13,19,21,22,23\}$ & $\{2,3,4,6,7,8,11,14,15,16,17,18,20,24\}$ & \texttt{011101110010\,011111010001} \\
$S_{8}$ & $\{2,5,11,13\}$ & $\{1,3,4,5,12,13,15,16,19,23\}$ & $\{2,6,7,8,9,10,11,14,17,18,20,21,22,24\}$ & \texttt{010001111110\,010011011101} \\
$S_{9}$ & $\{2,6,9,16\}$ & $\{1,5,9,10,12,15,17,21,22,24\}$ & $\{2,3,4,6,7,8,11,13,14,16,18,19,20,23\}$ & \texttt{011101110010\,110101110010} \\
$S_{10}$ & $\{2,6,10,15\}$ & $\{1,5,7,8,12,15,17,19,20,24\}$ & $\{2,3,4,6,9,10,11,13,14,16,18,21,22,23\}$ & \texttt{011101001110\,110101001110} \\
$S_{11}$ & $\{2,8,10,15\}$ & $\{1,5,7,8,12,14,18,19,20,22\}$ & $\{2,3,4,6,9,10,11,13,15,16,17,21,23,24\}$ & \texttt{011101001110\,101110001011} \\
$S_{12}$ & $\{2,8,11,13\}$ & $\{1,3,4,5,12,14,15,16,18,22\}$ & $\{2,6,7,8,9,10,11,13,17,19,20,21,23,24\}$ & \texttt{010001111110\,100010111011} \\
$S_{13}$ & $\{3,5,10,14\}$ & $\{2,5,6,8,9,13,17,18,19,23\}$ & $\{1,3,4,7,10,11,12,14,15,16,20,21,22,24\}$ & \texttt{101100100111\,011100011101} \\
$S_{14}$ & $\{3,5,11,13\}$ & $\{2,3,4,8,9,13,15,16,19,23\}$ & $\{1,5,6,7,10,11,12,14,17,18,20,21,22,24\}$ & \texttt{100011100111\,010011011101} \\
$S_{15}$ & $\{3,7,9,17\}$ & $\{2,8,9,11,12,16,20,21,23,24\}$ & $\{1,3,4,5,6,7,10,13,14,15,17,18,19,22\}$ & \texttt{101111100100\,111011100100} \\
$S_{16}$ & $\{3,7,10,14\}$ & $\{2,5,6,8,9,16,17,18,20,21\}$ & $\{1,3,4,7,10,11,12,13,14,15,19,22,23,24\}$ & \texttt{101100100111\,111000100111} \\
$S_{17}$ & $\{3,8,9,17\}$ & $\{2,8,9,11,12,14,18,22,23,24\}$ & $\{1,3,4,5,6,7,10,13,15,16,17,19,20,21\}$ & \texttt{101111100100\,101110111000} \\
$S_{18}$ & $\{3,8,11,13\}$ & $\{2,3,4,8,9,14,15,16,18,22\}$ & $\{1,5,6,7,10,11,12,13,17,19,20,21,23,24\}$ & \texttt{100011100111\,100010111011} \\
$S_{19}$ & $\{4,6,10,15\}$ & $\{4,6,7,8,10,15,17,19,20,24\}$ & $\{1,2,3,5,9,11,12,13,14,16,18,21,22,23\}$ & \texttt{111010001011\,110101001110} \\
$S_{20}$ & $\{4,6,11,12\}$ & $\{1,2,4,6,10,13,14,15,17,24\}$ & $\{3,5,7,8,9,11,12,16,18,19,20,21,22,23\}$ & \texttt{001010111011\,000101111110} \\
$S_{21}$ & $\{4,7,9,17\}$ & $\{4,6,10,11,12,16,20,21,23,24\}$ & $\{1,2,3,5,7,8,9,13,14,15,17,18,19,22\}$ & \texttt{111010111000\,111011100100} \\
$S_{22}$ & $\{4,7,11,12\}$ & $\{1,2,4,6,10,13,14,16,20,21\}$ & $\{3,5,7,8,9,11,12,15,17,18,19,22,23,24\}$ & \texttt{001010111011\,001011100111} \\
$S_{23}$ & $\{4,8,9,17\}$ & $\{4,6,10,11,12,14,18,22,23,24\}$ & $\{1,2,3,5,7,8,9,13,15,16,17,19,20,21\}$ & \texttt{111010111000\,101110111000} \\
$S_{24}$ & $\{4,8,10,15\}$ & $\{4,6,7,8,10,14,18,19,20,22\}$ & $\{1,2,3,5,9,11,12,13,15,16,17,21,23,24\}$ & \texttt{111010001011\,101110001011}
\end{tabular}
\end{ruledtabular}
\end{table*}

\begin{table*}[t]
\caption{Vertex state profiles and weights for all $65$ vertices of $H_{65}$, showing for each core vertex $i\in\{1,\ldots,17\}$ and each completion pair $(w_{m,1},w_{m,2})=(2m+16,2m+17)$ ($m=1,\ldots,24$) the number of active branches $W_{24}$, the weight $W_{48}=\sum_{\alpha,\varepsilon}s_{\alpha,\varepsilon}(i)$ across the $48$ canonical states, the weight $W_{393{,}216}=\sum_{s\in\mathcal{S}(H_{65})}s(i)$ across all $393{,}216$ states, the active branch indices $\alpha\in\{1,\ldots,24\}$ (equal to $Z(e_m)$ for $(w_{m,1},w_{m,2})$), and the $24$-bit branch profile across $S_1,\ldots,S_{24}$ (in which each \texttt{1} expands to \texttt{11} for core vertices $i$, to \texttt{10} for $w_{m,1}$, and to \texttt{01} for $w_{m,2}$ across the $48$ canonical states).}\label{tab:app-h65-vertex-profiles}
\scriptsize
\begin{ruledtabular}
\begin{tabular}{lrrrlc}
Vertex / completion pair & $W_{24}$ & $W_{48}$ & $W_{393{,}216}$ & Active branches $\alpha\in\{1,\ldots,24\}$ & $24$-bit branch profile $(1\ldots 8\,|\,9\ldots 16\,|\,17\ldots 24)$ \\ \hline
Core $1$ & $6$ & $12$ & $98{,}304$ & $\{1,2,3,4,5,6\}$ & \texttt{11111100\,00000000\,00000000} \\
Core $2$ & $6$ & $12$ & $98{,}304$ & $\{7,8,9,10,11,12\}$ & \texttt{00000011\,11110000\,00000000} \\
Core $3$ & $6$ & $12$ & $98{,}304$ & $\{13,14,15,16,17,18\}$ & \texttt{00000000\,00001111\,11000000} \\
Core $4$ & $6$ & $12$ & $98{,}304$ & $\{19,20,21,22,23,24\}$ & \texttt{00000000\,00000000\,00111111} \\
Core $5$ & $6$ & $12$ & $98{,}304$ & $\{1,2,7,8,13,14\}$ & \texttt{11000011\,00001100\,00000000} \\
Core $6$ & $6$ & $12$ & $98{,}304$ & $\{3,4,9,10,19,20\}$ & \texttt{00110000\,11000000\,00110000} \\
Core $7$ & $6$ & $12$ & $98{,}304$ & $\{5,6,15,16,21,22\}$ & \texttt{00001100\,00000011\,00001100} \\
Core $8$ & $6$ & $12$ & $98{,}304$ & $\{11,12,17,18,23,24\}$ & \texttt{00000000\,00110000\,11000011} \\
Core $9$ & $8$ & $16$ & $131{,}072$ & $\{1,3,7,9,15,17,21,23\}$ & \texttt{10100010\,10000010\,10001010} \\
Core $10$ & $8$ & $16$ & $131{,}072$ & $\{2,5,10,11,13,16,19,24\}$ & \texttt{01001000\,01101001\,00100001} \\
Core $11$ & $8$ & $16$ & $131{,}072$ & $\{4,6,8,12,14,18,20,22\}$ & \texttt{00010101\,00010100\,01010100} \\
Core $12$ & $4$ & $8$ & $65{,}536$ & $\{4,6,20,22\}$ & \texttt{00010100\,00000000\,00010100} \\
Core $13$ & $4$ & $8$ & $65{,}536$ & $\{8,12,14,18\}$ & \texttt{00000001\,00010100\,01000000} \\
Core $14$ & $4$ & $8$ & $65{,}536$ & $\{2,5,13,16\}$ & \texttt{01001000\,00001001\,00000000} \\
Core $15$ & $4$ & $8$ & $65{,}536$ & $\{10,11,19,24\}$ & \texttt{00000000\,01100000\,00100001} \\
Core $16$ & $4$ & $8$ & $65{,}536$ & $\{1,3,7,9\}$ & \texttt{10100010\,10000000\,00000000} \\
Core $17$ & $4$ & $8$ & $65{,}536$ & $\{15,17,21,23\}$ & \texttt{00000000\,00000010\,10001010} \\ \hline
$E_{1}:\,(18,19)$ & $14$ & $14$ & $114{,}688$ & $\{1,2,3,5,13,14,15,16,17,18,19,21,23,24\}$ & \texttt{11101000\,00001111\,11101011} \\
$E_{2}:\,(20,21)$ & $14$ & $14$ & $114{,}688$ & $\{1,2,3,5,7,8,9,10,11,12,19,21,23,24\}$ & \texttt{11101011\,11110000\,00101011} \\
$E_{3}:\,(22,23)$ & $14$ & $14$ & $114{,}688$ & $\{7,9,10,11,13,15,16,17,19,20,21,22,23,24\}$ & \texttt{00000010\,11101011\,10111111} \\
$E_{4}:\,(24,25)$ & $14$ & $14$ & $114{,}688$ & $\{1,2,3,4,5,6,7,9,10,11,13,15,16,17\}$ & \texttt{11111110\,11101011\,10000000} \\
$E_{5}:\,(26,27)$ & $14$ & $14$ & $114{,}688$ & $\{1,3,4,6,14,15,17,18,19,20,21,22,23,24\}$ & \texttt{10110100\,00000110\,11111111} \\
$E_{6}:\,(28,29)$ & $14$ & $14$ & $114{,}688$ & $\{1,3,4,6,7,8,9,10,11,12,14,15,17,18\}$ & \texttt{10110111\,11110110\,11000000} \\
$E_{7}:\,(30,31)$ & $14$ & $14$ & $114{,}688$ & $\{7,8,9,12,13,14,15,16,17,18,20,21,22,23\}$ & \texttt{00000011\,10011111\,11011110} \\
$E_{8}:\,(32,33)$ & $14$ & $14$ & $114{,}688$ & $\{1,2,3,4,5,6,7,8,9,12,20,21,22,23\}$ & \texttt{11111111\,10010000\,00011110} \\
$E_{9}:\,(34,35)$ & $14$ & $14$ & $114{,}688$ & $\{2,4,5,6,8,10,11,12,19,20,21,22,23,24\}$ & \texttt{01011101\,01110000\,00111111} \\
$E_{10}:\,(36,37)$ & $14$ & $14$ & $114{,}688$ & $\{2,4,5,6,8,10,11,12,13,14,15,16,17,18\}$ & \texttt{01011101\,01111111\,11000000} \\
$E_{11}:\,(38,39)$ & $14$ & $14$ & $114{,}688$ & $\{7,8,9,10,11,12,13,14,16,18,19,20,22,24\}$ & \texttt{00000011\,11111101\,01110101} \\
$E_{12}:\,(40,41)$ & $14$ & $14$ & $114{,}688$ & $\{1,2,3,4,5,6,13,14,16,18,19,20,22,24\}$ & \texttt{11111100\,00001101\,01110101} \\
$E_{13}:\,(42,43)$ & $14$ & $14$ & $114{,}688$ & $\{3,5,9,10,11,12,15,16,17,18,19,21,23,24\}$ & \texttt{00101000\,11110011\,11101011} \\
$E_{14}:\,(44,45)$ & $14$ & $14$ & $114{,}688$ & $\{1,2,3,5,7,8,9,10,13,14,15,16,19,21\}$ & \texttt{11101011\,11001111\,00101000} \\
$E_{15}:\,(46,47)$ & $14$ & $14$ & $114{,}688$ & $\{1,2,5,6,7,11,13,15,16,17,21,22,23,24\}$ & \texttt{11001110\,00101011\,10001111} \\
$E_{16}:\,(48,49)$ & $14$ & $14$ & $114{,}688$ & $\{1,2,3,4,7,9,10,11,13,17,19,20,23,24\}$ & \texttt{11110010\,11101000\,10110011} \\
$E_{17}:\,(50,51)$ & $14$ & $14$ & $114{,}688$ & $\{1,6,7,8,11,12,14,15,17,18,21,22,23,24\}$ & \texttt{10000111\,00110110\,11001111} \\
$E_{18}:\,(52,53)$ & $14$ & $14$ & $114{,}688$ & $\{1,3,4,6,7,8,9,10,14,15,19,20,21,22\}$ & \texttt{10110111\,11000110\,00111100} \\
$E_{19}:\,(54,55)$ & $14$ & $14$ & $114{,}688$ & $\{3,4,5,6,9,12,15,16,17,18,20,21,22,23\}$ & \texttt{00111100\,10010011\,11011110} \\
$E_{20}:\,(56,57)$ & $14$ & $14$ & $114{,}688$ & $\{1,2,3,4,7,8,9,12,13,14,17,18,20,23\}$ & \texttt{11110011\,10011100\,11010010} \\
$E_{21}:\,(58,59)$ & $14$ & $14$ & $114{,}688$ & $\{2,4,8,10,11,12,13,14,17,18,19,20,23,24\}$ & \texttt{01010001\,01111100\,11110011} \\
$E_{22}:\,(60,61)$ & $14$ & $14$ & $114{,}688$ & $\{2,4,5,6,8,10,13,14,15,16,19,20,21,22\}$ & \texttt{01011101\,01001111\,00111100} \\
$E_{23}:\,(62,63)$ & $14$ & $14$ & $114{,}688$ & $\{3,4,5,6,9,10,11,12,16,18,19,20,22,24\}$ & \texttt{00111100\,11110001\,01110101} \\
$E_{24}:\,(64,65)$ & $14$ & $14$ & $114{,}688$ & $\{1,2,5,6,7,8,11,12,13,14,16,18,22,24\}$ & \texttt{11001111\,00111101\,01000101}
\end{tabular}
\end{ruledtabular}
\end{table*}



\subsection{Yu--Oh-type inequalities and the limits of state independence}\label{app:state-independence}

The state tables yield a direct argument similar to Yu--Oh's~\cite{Yu-2012}, provided that its dimensional scope is specified. For either configuration, write $\Pi_i=v_iv_i^T/\|v_i\|^2$ for its normalized ray projector, and let $p_i(\rho)=\operatorname{tr}(\rho\Pi_i)$ for a density operator $\rho\ge0$ with $\operatorname{tr}\rho=1$. A noncontextual probability assignment here is a convex mixture of the two-valued states respecting every listed context. An inequality is \emph{state-independent in dimension four} only if the same fixed projectors violate it for every four-dimensional density operator. This requirement is stronger than a chromatic or fractional-chromatic gap~\cite{cabello2015sic}.

\paragraph*{The Yu--Oh sector and the apex suspension.}
For the warm-up $H_{\YO}^{(4)}$, Eq.~\eqref{eq:hyo4-triad-disjoint} shows that at most one of vertices $10,11,12,13$ is assigned $1$. The apex state assigns $1$ only to vertex $26$, whereas all other states vanish at $26$. Thus every state, and every convex mixture, satisfies the tight inequality
\begin{equation}\label{eq:hyo4-contextuality-bound}
\mathcal{Y}=p_{26}+\sum_{i=10}^{13}p_i\le1.
\end{equation}
The four normalized triad projectors satisfy
\begin{equation}\label{eq:hyo4-contextuality-operator}
\sum_{i=10}^{13}\Pi_i=\frac43(I_4-\Pi_{26}),
\end{equation}
because their outer products have diagonal sum $4/3$ on the first three coordinates and cancel off-diagonally. Consequently, quantum theory gives
\begin{equation}\label{eq:hyo4-contextuality-gap}
\mathcal{Y}(\rho)=1+\frac{1-p_{26}(\rho)}{3}.
\end{equation}
For every state supported on the three-dimensional subspace orthogonal to the apex, this is the Yu--Oh value $4/3>1$, independent of the preparation within that sector. More generally, Eq.~\eqref{eq:hyo4-contextuality-bound} is violated whenever $p_{26}(\rho)<1$. The sole exceptional density operator is $\rho=\Pi_{26}$, whose probabilities reproduce the apex state exactly. Hence the suspension inherits the Yu--Oh argument on its three-dimensional sector, but is not a state-independent contextuality set in the full four-dimensional space. The apex obstruction is the same mechanism discussed in Ref.~\cite[Theorem~1]{cabello2015sic}.

\paragraph*{A contextuality witness for the linear hypergraph.}
For $H_{65}$, the independence bound gives
\begin{equation}\label{eq:h65-contextuality-bound}
\mathcal{C}=\sum_{i=1}^{17}p_i\le4.
\end{equation}
In fact, equality holds for every convex mixture of the $393{,}216$ hypergraph states, since each has exactly four active core vertices. With $q=(1,1,1,-1)^T$ and $Q=qq^T/4=\Pi_q$, Eq.~\eqref{eq:core-projector-sum} gives
\begin{equation}\label{eq:h65-contextuality-gap}
\mathcal{C}(\rho)=4+\operatorname{tr}(\rho Q).
\end{equation}
This violates Eq.~\eqref{eq:h65-contextuality-bound} for every state with nonzero weight along $q$, including the maximally mixed state $I_4/4$, for which the value is $17/4$. It does not yet establish a violation for states supported on $q^\perp$.

\paragraph*{An explicit noncontextual quantum state for all 65 rays.}
The failure of strict violation on $q^\perp$ cannot be repaired merely by choosing another inequality for the same measurements. To see this, consider
\begin{equation}\label{eq:h65-noncontextual-density}
\rho_\star=\frac{I_4-Q}{3}.
\end{equation}
Its core probabilities are
\begin{equation}\label{eq:h65-noncontextual-core}
p_i(\rho_\star)=
\begin{cases}
1/4,&1\le i\le8,\\
1/3,&9\le i\le11,\\
1/6,&12\le i\le17.
\end{cases}
\end{equation}
These coincide with the uniform mixture of the $24$ core supports $S_\alpha$, as follows from the core branch weights in Table~\ref{tab:app-h65-vertex-profiles}. It remains to reproduce the $48$ completion probabilities, rather than only the core.

For each completion context $E_m$, put $r_{m,j}=p_{w_{m,j}}(\rho_\star)$. Its two core endpoints have probabilities $1/4$ and $1/6$, so
\begin{equation}\label{eq:h65-noncontextual-completion}
r_{m,1}+r_{m,2}=\frac7{12},\qquad
\tau_m=\frac{12}{7}r_{m,2}\in[0,1].
\end{equation}
Choose $\alpha$ uniformly from $\{1,\ldots,24\}$ and independently choose $\lambda$ uniformly from $[0,1]$. Use $S_\alpha$ on the core. On each free context $m\in\mathcal{F}_\alpha$, set $\varepsilon_m=1$ if $\lambda<\tau_m$ and $\varepsilon_m=0$ otherwise, and use Eq.~\eqref{eq:full-state-param}; on hit contexts both completion values remain zero. Every resulting assignment is therefore one of the already enumerated two-valued states. Since precisely $14$ branches miss any given bridge edge, the completion marginals are
\begin{equation}\label{eq:h65-noncontextual-marginals}
\begin{aligned}
\mathbb{E}[s(w_{m,2})]&=\frac{14}{24}\tau_m=r_{m,2},\\
\mathbb{E}[s(w_{m,1})]&=\frac{14}{24}(1-\tau_m)=r_{m,1}.
\end{aligned}
\end{equation}
Together with Eq.~\eqref{eq:h65-noncontextual-core}, this constructs a noncontextual model for all $65$ quantum probabilities of $\rho_\star$. The probabilities for a complete context determine its full outcome distribution, so the model also reproduces every listed measurement context. The same holds for compatible subsets of the named ray projectors, since distinct orthogonal rank-one projectors have zero joint products.

For the coordinates given here, $\tau_m$ takes only the four values $4/9$, $5/11$, $6/11$, and $5/9$. Partitioning $[0,1]$ at these thresholds expresses the model as a mixture of at most $5\times24=120$ of the states parametrized in Appendix~\ref{app:states-h65}. Thus $H_{65}$, with this fixed realization, also admits no fully state-independent contextuality inequality for the specified ray measurements. A preparation-dependent rotation of the projectors could make $\operatorname{tr}(\rho Q)>0$ for any given preparation, but would change the measurements with the preparation and would not be a state-independent test using a fixed set~\cite{cabello2015sic}.

\paragraph*{A fixed enlargement giving a state-independent argument.}
There is nevertheless a simple way to remove the exceptional subspace by enlarging the measurement family. Let
\begin{equation}\label{eq:fixed-core-rotations}
\begin{aligned}
D_0&=I_4,&D_1&=\operatorname{diag}(1,1,-1,-1),\\
D_2&=\operatorname{diag}(1,-1,1,-1),&
D_3&=\operatorname{diag}(1,-1,-1,1),
\end{aligned}
\end{equation}
and define $\Pi_i^{(a)}=D_a\Pi_iD_a^T$ for $1\le i\le17$ and $0\le a\le3$. The four vectors $D_aq/2$ form an orthonormal basis. Applying Eq.~\eqref{eq:core-projector-sum} to each copy therefore gives
\begin{equation}\label{eq:symmetrized-core-operator}
\sum_{a=0}^{3}\sum_{i=1}^{17}\Pi_i^{(a)}
=16I_4+\sum_{a=0}^{3}D_aQD_a^T=17I_4.
\end{equation}
Any deterministic assignment respecting pairwise orthogonality restricts to an independent set of size at most $4$ on each rotated core. Hence every convex mixture of such assignments satisfies
\begin{equation}\label{eq:symmetrized-core-inequality}
\sum_{a=0}^{3}\sum_{i=1}^{17}p_i^{(a)}\le16,
\qquad\text{whereas quantum theory gives }17
\end{equation}
for every density operator. Repeated projectors in this sum are the same measurement event and are counted with their multiplicity. The classical bound $16$ suffices for the argument; it is not claimed to be tight. All four rotations are fixed independently of the preparation.

This construction introduces seven additional core rays: its distinct rays are exactly the $24$ Peres rays~\cite{peres-91}, namely the four coordinate rays, the twelve rays $e_j\pm e_k$ ($1\le j<k\le4$), and the eight rays $(1,\pm1,\pm1,\pm1)^T$. Including all their orthogonal bases recovers the Kochen--Specker obstruction of that configuration. Thus Eq.~\eqref{eq:symmetrized-core-inequality} supplies a state-independent argument for an enlarged family, not for the original $65$ fixed rays with their separating two-valued states.

Both original hypergraphs consequently support Kochen--Specker-type probability inequalities while retaining their two-valued states. The warm-up supplies the state-independent Yu--Oh argument on its three-dimensional sector; a fully four-dimensional state-independent argument requires an enlargement such as the one above.

\section{Reproducible completion rule and exact verification}\label{app:verification}

The coordinates in Eq.~\eqref{eq:h65-coordinates} can be reproduced by the following finite selection rule. Normalize an integer vector by dividing by the greatest common divisor of its entries and choosing the sign so that its first nonzero entry is positive. Form the set of normalized nonzero vectors obtained from $\{-3,-2,\ldots,3\}^4$, ordered by squared norm, then maximum absolute entry, then lexicographically. Process bridge edges in the order $e_{13},\ldots,e_{24},e_1,\ldots,e_{12}$. For each pair, choose $\eta_A$ to be the first vector in this ordered pool orthogonal to its endpoints, and $\eta_B$ the first such vector also orthogonal to $\eta_A$.

For all $(a,b)\in\{-7,-6,\ldots,7\}^2\setminus\{(0,0)\}$, form the pair in Eq.~\eqref{eq:plane-param} and normalize both vectors. Discard pairs that repeat an already placed ray or have an orthogonality to any already placed vertex other than the two endpoints. Among the remaining ordered pairs choose the lexicographically smallest score
\begin{equation}\label{eq:completion-score}
\left(\max_{j,k}|(w_{m,j})_k|,\
\|w_{m,1}\|^2+\|w_{m,2}\|^2,\ w_{m,1},w_{m,2}\right).
\end{equation}
In Eq.~\eqref{eq:completion-score}, the order on the last two entries is the usual lexicographic order on their integer coordinates. This rule yields exactly the $48$ completion vectors in Table~\ref{tab:app-h65-vectors}. It is a recipe for this specific example; the general existence assertion is established independently by Lemma~\ref{lem:plane-completion}.

The supplementary Python~3.10 script invoked below parses this TeX source and reproduces the completion rule, coordinate and state-profile tables, all $393{,}216$ state extensions, and the fractional-coloring certificate. It checks every orthogonality, the maximal cliques, atomic-order failure, the forced hub-completion zero, and the analogous warm-up data. It uses only the standard library and writes a machine-readable report. Keep the script, source, and master bibliography together and run
\begin{verbatim}
python 2026-chroma-dim4-example-verify.py
\end{verbatim}
The source, script, and report accompany the manuscript.

\begin{acknowledgments}
This research was funded in part by the Austrian Science Fund (FWF), Grant DOI \href{https://doi.org/10.55776/PIN5424624}{10.55776/PIN5424624}, and the Czech Science Foundation (GA\v{C}R), Grant No.~25-20013L.
The author acknowledges FWF and TU Wien Bibliothek for financial support through its Open Access Funding Programme, as well as support from Google via Gemini 4 Argon High and contributions from SOL 6.2 xhigh and Sonnet 5.5 high.
\end{acknowledgments}

\bibliography{svozil}

\end{document}
